
\documentclass{article}
\usepackage{colortbl}
\usepackage{makecell}
\usepackage{multirow}
\usepackage{supertabular}

\begin{document}

\newcounter{utterance}

\centering \large Interaction Transcript for game `clean\_up', experiment `1\_medium\_7obj\_nl', episode 4 with gemma{-}4{-}26b{-}a4b{-}it.
\vspace{24pt}

{ \footnotesize  \setcounter{utterance}{1}
\setlength{\tabcolsep}{0pt}
\begin{supertabular}{c@{$\;$}|p{.15\linewidth}@{}p{.15\linewidth}p{.15\linewidth}p{.15\linewidth}p{.15\linewidth}p{.15\linewidth}}
    \# & $\;$A & \multicolumn{4}{c}{Game Master} & $\;\:$B\\
    \hline

    \theutterance \stepcounter{utterance}  
    & & \multicolumn{4}{p{0.6\linewidth}}{
        \cellcolor[rgb]{0.9,0.9,0.9}{
            \makecell[{{p{\linewidth}}}]{
                \texttt{\tiny{[P1$\langle$GM]}}
                \texttt{Ik ben uw spelleider, en u speelt een samenwerkingsspel met het volgende raster als speelbord:} \\
\\ 
\texttt{```} \\
\texttt{ 1234567} \\
\texttt{╔═╤══╤══╗} \\
\texttt{║◌│◌◌│◌◌║ 1} \\
\texttt{╟─┤◌◌│◌S║ 2} \\
\texttt{║◌│TD├──╢ 3} \\
\texttt{╟─┘◌◌│O◌║ 4} \\
\texttt{║◌◌PM│◌A║ 5} \\
\texttt{║◌◌┌─┼──╢ 6} \\
\texttt{║◌◌│◌│◌◌║ 7} \\
\texttt{╚══╧═╧══╝} \\
\\ 
\texttt{```} \\
\\ 
\texttt{* De bovenrand toont x{-}coördinaten die naar rechts oplopen, en de rechterrand y{-}coördinaten die naar beneden oplopen.} \\
\texttt{* De volgende objecten zijn willekeurig op uw raster geplaatst: 'P', 'O', 'T', 'S', 'D', 'A', 'M'.} \\
\\ 
\\ 
\texttt{De andere speler ziet een variatie van het speelbord, waarbij de objecten op verschillende willekeurige locaties zijn geplaatst. U kunt het bord van de andere speler niet zien en zij kunnen het uwe niet zien.} \\
\\ 
\\ 
\texttt{**Doel:**} \\
\\ 
\texttt{Beide spelers moeten de objecten op hun respectievelijke achtergrond zo verplaatsen dat identieke objecten op dezelfde coördinaten terechtkomen. U moet communiceren met de andere speler om een gemeenschappelijke doelconfiguratie overeen te komen.} \\
\\ 
\\ 
\texttt{**Regels:**} \\
\\ 
\texttt{* In elke beurt kunt u precies één van de volgende twee commando's sturen:} \\
\texttt{ 1. `ZEG: <BERICHT>`: om een bericht (alles tot de volgende regeleinde) naar de andere speler te sturen. Ik zal het doorsturen naar uw partner.} \\
\texttt{2. `BEWEEG: <OBJECT>, (<X>, <Y>)`: om een object naar een nieuwe positie te verplaatsen, waarbij `<X>` de kolom is en `<Y>` de rij. Ik zal u laten weten of uw zet succesvol was of niet.} \\
\texttt{* Als u zich niet aan het formaat houdt, of meerdere commando's tegelijk verstuurt, moet ik u straffen.} \\
\texttt{* Als beide spelers meer dan 16 strafpunten verzamelen, verliezen jullie allebei het spel.} \\
\texttt{* Het is van vitaal belang dat u met de andere speler communiceert over uw doelstatus! De *enige* manier waarop u uw strategie aan de andere speler kunt overbrengen is door gebruik te maken van het `ZEG: <BERICHT>` commando!} \\
\\ 
\\ 
\texttt{**Objecten verplaatsen**} \\
\\ 
\texttt{* U kunt alleen objecten verplaatsen naar cellen binnen de grenzen van het raster. De doelcel moet leeg zijn, d.w.z. deze mag alleen het symbool '◌' bevatten.} \\
\texttt{* Als u probeert een object te verplaatsen naar een plek die niet leeg is, of als u probeert het buiten het raster te verplaatsen, moet ik u straffen. U krijgt een nieuwe poging.} \\
\texttt{* Voordat u een zet doet, moet u controleren of de doelplek leeg is, en geen letter, kader of lijn bevat!} \\
\\ 
\\ 
\texttt{**Einde van het spel**} \\
\\ 
\texttt{Als u denkt dat u het doel bereikt hebt om alle objecten uit te lijnen, kunt u de andere speler vragen het spel te beëindigen door `ZEG: klaar?` te sturen. Als de andere speler u vraagt om het spel te beëindigen en u antwoordt met `ZEG: klaar!`, zal het spel eindigen.} \\
\\ 
\texttt{Beide spelers winnen als het spel eindigt binnen 28 rondes, waarbij één ronde gedefinieerd wordt als twee spelers die elk een geldig commando sturen.} \\
\\ 
\texttt{**Score:**} \\
\\ 
\texttt{Hoe dichter de identieke objecten bij elkaar liggen op beide speelborden, hoe meer punten u krijgt. Strafpunten verminderen uw punten. Kunt u het record verslaan?} \\
\\ 
\\ 
\texttt{Stuur een bericht naar de andere speler om het spel te starten!} \\
            }
        }
    }
    & \\ \\

    \theutterance \stepcounter{utterance}  
    & \multicolumn{4}{p{0.6\linewidth}}{
        \cellcolor[rgb]{0.9,0.9,0.9}{
            \makecell[{{p{\linewidth}}}]{
                \texttt{\tiny{[P1$\rangle$GM]}}
                \texttt{ZEG: Hallo partner! Laten we beginnen met het afstemmen van onze strategie. Ik zie de volgende objecten op mijn bord: S op (7,2), T op (2,3), D op (3,3), O op (6,4), P op (3,5), M op (4,5) en A op (7,5). Zullen we proberen om alle objecten in een compact blok in het midden of een specifieke vorm te plaatsen? Laat me weten welke objecten jij ziet en welk doelconfiguratie jij voorstelt!} \\
            }
        }
    }
    & & \\ \\

    \theutterance \stepcounter{utterance}  
    & & & \multicolumn{2}{p{0.3\linewidth}}{
        \cellcolor[rgb]{0.9,0.9,0.9}{
            \makecell[{{p{\linewidth}}}]{
                \texttt{\tiny{[GM$|$GM]}}
                \texttt{ZEG: Hallo partner! Laten we beginnen met het afstemmen van onze strategie. Ik zie de volgende objecten op mijn bord: S op (7,2), T op (2,3), D op (3,3), O op (6,4), P op (3,5), M op (4,5) en A op (7,5). Zullen we proberen om alle objecten in een compact blok in het midden of een specifieke vorm te plaatsen? Laat me weten welke objecten jij ziet en welk doelconfiguratie jij voorstelt!} \\
            }
        }
    }
    & & \\ \\

    \theutterance \stepcounter{utterance}  
    & & \multicolumn{4}{p{0.6\linewidth}}{
        \cellcolor[rgb]{0.9,0.9,0.9}{
            \makecell[{{p{\linewidth}}}]{
                \texttt{\tiny{[GM$\rangle$P2]}}
                \texttt{Ik ben uw spelleider, en u speelt een samenwerkingsspel met het volgende raster als speelbord:} \\
\\ 
\texttt{```} \\
\texttt{ 1234567} \\
\texttt{╔═╤══╤══╗} \\
\texttt{║◌│◌O│T◌║ 1} \\
\texttt{╟─┤◌◌│◌◌║ 2} \\
\texttt{║◌│◌◌├──╢ 3} \\
\texttt{╟─┘◌D│◌◌║ 4} \\
\texttt{║◌SA◌│◌◌║ 5} \\
\texttt{║PM┌─┼──╢ 6} \\
\texttt{║◌◌│◌│◌◌║ 7} \\
\texttt{╚══╧═╧══╝} \\
\\ 
\texttt{```} \\
\\ 
\texttt{* De bovenrand toont x{-}coördinaten die naar rechts oplopen, en de rechterrand y{-}coördinaten die naar beneden oplopen.} \\
\texttt{* De volgende objecten zijn willekeurig op uw raster geplaatst: 'P', 'O', 'T', 'S', 'D', 'A', 'M'.} \\
\\ 
\\ 
\texttt{De andere speler ziet een variatie van het speelbord, waarbij de objecten op verschillende willekeurige locaties zijn geplaatst. U kunt het bord van de andere speler niet zien en zij kunnen het uwe niet zien.} \\
\\ 
\\ 
\texttt{**Doel:**} \\
\\ 
\texttt{Beide spelers moeten de objecten op hun respectievelijke achtergrond zo verplaatsen dat identieke objecten op dezelfde coördinaten terechtkomen. U moet communiceren met de andere speler om een gemeenschappelijke doelconfiguratie overeen te komen.} \\
\\ 
\\ 
\texttt{**Regels:**} \\
\\ 
\texttt{* In elke beurt kunt u precies één van de volgende twee commando's sturen:} \\
\texttt{ 1. `ZEG: <BERICHT>`: om een bericht (alles tot de volgende regeleinde) naar de andere speler te sturen. Ik zal het doorsturen naar uw partner.} \\
\texttt{2. `BEWEEG: <OBJECT>, (<X>, <Y>)`: om een object naar een nieuwe positie te verplaatsen, waarbij `<X>` de kolom is en `<Y>` de rij. Ik zal u laten weten of uw zet succesvol was of niet.} \\
\texttt{* Als u zich niet aan het formaat houdt, of meerdere commando's tegelijk verstuurt, moet ik u straffen.} \\
\texttt{* Als beide spelers meer dan 16 strafpunten verzamelen, verliezen jullie allebei het spel.} \\
\texttt{* Het is van vitaal belang dat u met de andere speler communiceert over uw doelstatus! De *enige* manier waarop u uw strategie aan de andere speler kunt overbrengen is door gebruik te maken van het `ZEG: <BERICHT>` commando!} \\
\\ 
\\ 
\texttt{**Objecten verplaatsen**} \\
\\ 
\texttt{* U kunt alleen objecten verplaatsen naar cellen binnen de grenzen van het raster. De doelcel moet leeg zijn, d.w.z. deze mag alleen het symbool '◌' bevatten.} \\
\texttt{* Als u probeert een object te verplaatsen naar een plek die niet leeg is, of als u probeert het buiten het raster te verplaatsen, moet ik u straffen. U krijgt een nieuwe poging.} \\
\texttt{* Voordat u een zet doet, moet u controleren of de doelplek leeg is, en geen letter, kader of lijn bevat!} \\
\\ 
\\ 
\texttt{**Einde van het spel**} \\
\\ 
\texttt{Als u denkt dat u het doel bereikt hebt om alle objecten uit te lijnen, kunt u de andere speler vragen het spel te beëindigen door `ZEG: klaar?` te sturen. Als de andere speler u vraagt om het spel te beëindigen en u antwoordt met `ZEG: klaar!`, zal het spel eindigen.} \\
\\ 
\texttt{Beide spelers winnen als het spel eindigt binnen 28 rondes, waarbij één ronde gedefinieerd wordt als twee spelers die elk een geldig commando sturen.} \\
\\ 
\texttt{**Score:**} \\
\\ 
\texttt{Hoe dichter de identieke objecten bij elkaar liggen op beide speelborden, hoe meer punten u krijgt. Strafpunten verminderen uw punten. Kunt u het record verslaan?} \\
\\ 
\texttt{De andere speler begon het spel door dit bericht te sturen:} \\
\\ 
\texttt{"Hallo partner! Laten we beginnen met het afstemmen van onze strategie. Ik zie de volgende objecten op mijn bord: S op (7,2), T op (2,3), D op (3,3), O op (6,4), P op (3,5), M op (4,5) en A op (7,5). Zullen we proberen om alle objecten in een compact blok in het midden of een specifieke vorm te plaatsen? Laat me weten welke objecten jij ziet en welk doelconfiguratie jij voorstelt!"} \\
\\ 
\texttt{Wat is uw eerste opdracht?} \\
            }
        }
    }
    & \\ \\

    \theutterance \stepcounter{utterance}  
    & & & \multicolumn{4}{p{0.6\linewidth}}{
        \cellcolor[rgb]{0.9,0.9,0.9}{
            \makecell[{{p{\linewidth}}}]{
                \texttt{\tiny{[GM$\langle$P2]}}
                \texttt{ZEG: Hallo partner! Ik heb de volgende objecten op mijn bord: O op (4,1), T op (6,1), D op (4,4), S op (2,5), A op (3,5), P op (1,6) en M op (1,6) {-} wacht, P en M lijken op (1,6) en (2,6) te staan volgens het raster. Laten we proberen alle objecten in een compact blok in het midden te plaatsen, bijvoorbeeld in de zone rond (3,3) tot (5,5). Wat vind je van de volgende doelposities: P(3,3), M(4,3), S(3,4), A(4,4), D(3,5), O(4,5) en T(5,5)?} \\
            }
        }
    }
     \\ \\

    \theutterance \stepcounter{utterance}  
    & & & \multicolumn{2}{p{0.3\linewidth}}{
        \cellcolor[rgb]{0.9,0.9,0.9}{
            \makecell[{{p{\linewidth}}}]{
                \texttt{\tiny{[GM$|$GM]}}
                \texttt{ZEG: Hallo partner! Ik heb de volgende objecten op mijn bord: O op (4,1), T op (6,1), D op (4,4), S op (2,5), A op (3,5), P op (1,6) en M op (1,6) {-} wacht, P en M lijken op (1,6) en (2,6) te staan volgens het raster. Laten we proberen alle objecten in een compact blok in het midden te plaatsen, bijvoorbeeld in de zone rond (3,3) tot (5,5). Wat vind je van de volgende doelposities: P(3,3), M(4,3), S(3,4), A(4,4), D(3,5), O(4,5) en T(5,5)?} \\
            }
        }
    }
    & & \\ \\

    \theutterance \stepcounter{utterance}  
    & & \multicolumn{4}{p{0.6\linewidth}}{
        \cellcolor[rgb]{0.9,0.9,0.9}{
            \makecell[{{p{\linewidth}}}]{
                \texttt{\tiny{[P1$\langle$GM]}}
                \texttt{Je bericht is doorgestuurd naar de andere speler.} \\
\\ 
\\ 
\texttt{Je speelt momenteel ronde 1 van maximaal 28.} \\
\texttt{Jullie hebben gezamenlijk 0 van 16 straffen verzameld.} \\
\texttt{De andere speler stuurde dit bericht:} \\
\\ 
\\ 
\texttt{"Hallo partner! Ik heb de volgende objecten op mijn bord: O op (4,1), T op (6,1), D op (4,4), S op (2,5), A op (3,5), P op (1,6) en M op (1,6) {-} wacht, P en M lijken op (1,6) en (2,6) te staan volgens het raster. Laten we proberen alle objecten in een compact blok in het midden te plaatsen, bijvoorbeeld in de zone rond (3,3) tot (5,5). Wat vind je van de volgende doelposities: P(3,3), M(4,3), S(3,4), A(4,4), D(3,5), O(4,5) en T(5,5)?"} \\
\\ 
\\ 
\texttt{Wat is je volgende opdracht?} \\
            }
        }
    }
    & \\ \\

    \theutterance \stepcounter{utterance}  
    & \multicolumn{4}{p{0.6\linewidth}}{
        \cellcolor[rgb]{0.9,0.9,0.9}{
            \makecell[{{p{\linewidth}}}]{
                \texttt{\tiny{[P1$\rangle$GM]}}
                \texttt{ZEG: Dat is een goed plan! Laten we die doelposities aanhouden: P(3,3), M(4,3), S(3,4), A(4,4), D(3,5), O(4,5) en T(5,5). Ik ga beginnen met het verplaatsen van mijn objecten naar deze posities. Ik begin met het verplaatsen van mijn T naar (5,5).} \\
            }
        }
    }
    & & \\ \\

    \theutterance \stepcounter{utterance}  
    & & & \multicolumn{2}{p{0.3\linewidth}}{
        \cellcolor[rgb]{0.9,0.9,0.9}{
            \makecell[{{p{\linewidth}}}]{
                \texttt{\tiny{[GM$|$GM]}}
                \texttt{ZEG: Dat is een goed plan! Laten we die doelposities aanhouden: P(3,3), M(4,3), S(3,4), A(4,4), D(3,5), O(4,5) en T(5,5). Ik ga beginnen met het verplaatsen van mijn objecten naar deze posities. Ik begin met het verplaatsen van mijn T naar (5,5).} \\
            }
        }
    }
    & & \\ \\

    \theutterance \stepcounter{utterance}  
    & & \multicolumn{4}{p{0.6\linewidth}}{
        \cellcolor[rgb]{0.9,0.9,0.9}{
            \makecell[{{p{\linewidth}}}]{
                \texttt{\tiny{[GM$\rangle$P2]}}
                \texttt{Je bericht is doorgestuurd naar de andere speler.} \\
\\ 
\\ 
\texttt{Je speelt momenteel ronde 2 van maximaal 28.} \\
\texttt{Jullie hebben gezamenlijk 0 van 16 straffen verzameld.} \\
\texttt{De andere speler stuurde dit bericht:} \\
\\ 
\\ 
\texttt{"Dat is een goed plan! Laten we die doelposities aanhouden: P(3,3), M(4,3), S(3,4), A(4,4), D(3,5), O(4,5) en T(5,5). Ik ga beginnen met het verplaatsen van mijn objecten naar deze posities. Ik begin met het verplaatsen van mijn T naar (5,5)."} \\
\\ 
\\ 
\texttt{Wat is je volgende opdracht?} \\
            }
        }
    }
    & \\ \\

    \theutterance \stepcounter{utterance}  
    & & & \multicolumn{4}{p{0.6\linewidth}}{
        \cellcolor[rgb]{0.9,0.9,0.9}{
            \makecell[{{p{\linewidth}}}]{
                \texttt{\tiny{[GM$\langle$P2]}}
                \texttt{BEWEEG: T, (6,1)} \\
            }
        }
    }
     \\ \\

    \theutterance \stepcounter{utterance}  
    & & & \multicolumn{2}{p{0.3\linewidth}}{
        \cellcolor[rgb]{0.9,0.9,0.9}{
            \makecell[{{p{\linewidth}}}]{
                \texttt{\tiny{[GM$|$GM]}}
                \texttt{BEWEEG: T, (6,1)} \\
            }
        }
    }
    & & \\ \\

    \theutterance \stepcounter{utterance}  
    & & & \multicolumn{2}{p{0.3\linewidth}}{
        \cellcolor[rgb]{0.9,0.9,0.9}{
            \makecell[{{p{\linewidth}}}]{
                \texttt{\tiny{[GM$|$GM]}}
                \texttt{Straf: (6,1) is niet leeg, maar bevat 'T'.} \\
\\ 
\texttt{Jullie hebben gezamenlijk 1 van 16 straffen verzameld.} \\
\texttt{Probeer het opnieuw!} \\
            }
        }
    }
    & & \\ \\

    \theutterance \stepcounter{utterance}  
    & & \multicolumn{4}{p{0.6\linewidth}}{
        \cellcolor[rgb]{0.9,0.9,0.9}{
            \makecell[{{p{\linewidth}}}]{
                \texttt{\tiny{[GM$\rangle$P2]}}
                \texttt{Straf: (6,1) is niet leeg, maar bevat 'T'.} \\
\\ 
\texttt{Jullie hebben gezamenlijk 1 van 16 straffen verzameld.} \\
\texttt{Probeer het opnieuw!} \\
            }
        }
    }
    & \\ \\

    \theutterance \stepcounter{utterance}  
    & & & \multicolumn{4}{p{0.6\linewidth}}{
        \cellcolor[rgb]{0.9,0.9,0.9}{
            \makecell[{{p{\linewidth}}}]{
                \texttt{\tiny{[GM$\langle$P2]}}
                \texttt{BEWEEG: T, (5,5)} \\
            }
        }
    }
     \\ \\

    \theutterance \stepcounter{utterance}  
    & & & \multicolumn{2}{p{0.3\linewidth}}{
        \cellcolor[rgb]{0.9,0.9,0.9}{
            \makecell[{{p{\linewidth}}}]{
                \texttt{\tiny{[GM$|$GM]}}
                \texttt{BEWEEG: T, (5,5)} \\
            }
        }
    }
    & & \\ \\

    \theutterance \stepcounter{utterance}  
    & & & \multicolumn{2}{p{0.3\linewidth}}{
        \cellcolor[rgb]{0.9,0.9,0.9}{
            \makecell[{{p{\linewidth}}}]{
                \texttt{\tiny{[GM$|$GM]}}
                \texttt{Straf: (5,5) is niet leeg, maar bevat '│'.} \\
\\ 
\texttt{Jullie hebben gezamenlijk 2 van 16 straffen verzameld.} \\
\texttt{Probeer het opnieuw!} \\
            }
        }
    }
    & & \\ \\

    \theutterance \stepcounter{utterance}  
    & & \multicolumn{4}{p{0.6\linewidth}}{
        \cellcolor[rgb]{0.9,0.9,0.9}{
            \makecell[{{p{\linewidth}}}]{
                \texttt{\tiny{[GM$\rangle$P2]}}
                \texttt{Straf: (5,5) is niet leeg, maar bevat '│'.} \\
\\ 
\texttt{Jullie hebben gezamenlijk 2 van 16 straffen verzameld.} \\
\texttt{Probeer het opnieuw!} \\
            }
        }
    }
    & \\ \\

    \theutterance \stepcounter{utterance}  
    & & & \multicolumn{4}{p{0.6\linewidth}}{
        \cellcolor[rgb]{0.9,0.9,0.9}{
            \makecell[{{p{\linewidth}}}]{
                \texttt{\tiny{[GM$\langle$P2]}}
                \texttt{BEWEEG: T, (7,1)} \\
            }
        }
    }
     \\ \\

    \theutterance \stepcounter{utterance}  
    & & & \multicolumn{2}{p{0.3\linewidth}}{
        \cellcolor[rgb]{0.9,0.9,0.9}{
            \makecell[{{p{\linewidth}}}]{
                \texttt{\tiny{[GM$|$GM]}}
                \texttt{BEWEEG: T, (7,1)} \\
            }
        }
    }
    & & \\ \\

    \theutterance \stepcounter{utterance}  
    & & & \multicolumn{2}{p{0.3\linewidth}}{
        \cellcolor[rgb]{0.9,0.9,0.9}{
            \makecell[{{p{\linewidth}}}]{
                \texttt{\tiny{[GM$|$GM]}}
                \texttt{T moved to (7, 1)} \\
            }
        }
    }
    & & \\ \\

    \theutterance \stepcounter{utterance}  
    & & \multicolumn{4}{p{0.6\linewidth}}{
        \cellcolor[rgb]{0.9,0.9,0.9}{
            \makecell[{{p{\linewidth}}}]{
                \texttt{\tiny{[P1$\langle$GM]}}
                \texttt{Je bericht is doorgestuurd naar de andere speler.} \\
\\ 
\\ 
\texttt{Je speelt momenteel ronde 2 van maximaal 28.} \\
\texttt{Jullie hebben gezamenlijk 2 van 16 straffen verzameld.} \\
\texttt{De andere speler heeft een voorwerp op zijn rooster verplaatst.} \\
\\ 
\\ 
\texttt{Wat is je volgende opdracht?} \\
            }
        }
    }
    & \\ \\

    \theutterance \stepcounter{utterance}  
    & \multicolumn{4}{p{0.6\linewidth}}{
        \cellcolor[rgb]{0.9,0.9,0.9}{
            \makecell[{{p{\linewidth}}}]{
                \texttt{\tiny{[P1$\rangle$GM]}}
                \texttt{BEWEEG: T, (5,5)} \\
            }
        }
    }
    & & \\ \\

    \theutterance \stepcounter{utterance}  
    & & & \multicolumn{2}{p{0.3\linewidth}}{
        \cellcolor[rgb]{0.9,0.9,0.9}{
            \makecell[{{p{\linewidth}}}]{
                \texttt{\tiny{[GM$|$GM]}}
                \texttt{BEWEEG: T, (5,5)} \\
            }
        }
    }
    & & \\ \\

    \theutterance \stepcounter{utterance}  
    & & & \multicolumn{2}{p{0.3\linewidth}}{
        \cellcolor[rgb]{0.9,0.9,0.9}{
            \makecell[{{p{\linewidth}}}]{
                \texttt{\tiny{[GM$|$GM]}}
                \texttt{Straf: (5,5) is niet leeg, maar bevat '│'.} \\
\\ 
\texttt{Jullie hebben gezamenlijk 3 van 16 straffen verzameld.} \\
\texttt{Probeer het opnieuw!} \\
            }
        }
    }
    & & \\ \\

    \theutterance \stepcounter{utterance}  
    & & \multicolumn{4}{p{0.6\linewidth}}{
        \cellcolor[rgb]{0.9,0.9,0.9}{
            \makecell[{{p{\linewidth}}}]{
                \texttt{\tiny{[P1$\langle$GM]}}
                \texttt{Straf: (5,5) is niet leeg, maar bevat '│'.} \\
\\ 
\texttt{Jullie hebben gezamenlijk 3 van 16 straffen verzameld.} \\
\texttt{Probeer het opnieuw!} \\
            }
        }
    }
    & \\ \\

    \theutterance \stepcounter{utterance}  
    & \multicolumn{4}{p{0.6\linewidth}}{
        \cellcolor[rgb]{0.9,0.9,0.9}{
            \makecell[{{p{\linewidth}}}]{
                \texttt{\tiny{[P1$\rangle$GM]}}
                \texttt{BEWEEG: T, (5,4)} \\
            }
        }
    }
    & & \\ \\

    \theutterance \stepcounter{utterance}  
    & & & \multicolumn{2}{p{0.3\linewidth}}{
        \cellcolor[rgb]{0.9,0.9,0.9}{
            \makecell[{{p{\linewidth}}}]{
                \texttt{\tiny{[GM$|$GM]}}
                \texttt{BEWEEG: T, (5,4)} \\
            }
        }
    }
    & & \\ \\

    \theutterance \stepcounter{utterance}  
    & & & \multicolumn{2}{p{0.3\linewidth}}{
        \cellcolor[rgb]{0.9,0.9,0.9}{
            \makecell[{{p{\linewidth}}}]{
                \texttt{\tiny{[GM$|$GM]}}
                \texttt{Straf: (5,4) is niet leeg, maar bevat '│'.} \\
\\ 
\texttt{Jullie hebben gezamenlijk 4 van 16 straffen verzameld.} \\
\texttt{Probeer het opnieuw!} \\
            }
        }
    }
    & & \\ \\

    \theutterance \stepcounter{utterance}  
    & & \multicolumn{4}{p{0.6\linewidth}}{
        \cellcolor[rgb]{0.9,0.9,0.9}{
            \makecell[{{p{\linewidth}}}]{
                \texttt{\tiny{[P1$\langle$GM]}}
                \texttt{Straf: (5,4) is niet leeg, maar bevat '│'.} \\
\\ 
\texttt{Jullie hebben gezamenlijk 4 van 16 straffen verzameld.} \\
\texttt{Probeer het opnieuw!} \\
            }
        }
    }
    & \\ \\

    \theutterance \stepcounter{utterance}  
    & \multicolumn{4}{p{0.6\linewidth}}{
        \cellcolor[rgb]{0.9,0.9,0.9}{
            \makecell[{{p{\linewidth}}}]{
                \texttt{\tiny{[P1$\rangle$GM]}}
                \texttt{BEWEEG: T, (5,1)} \\
            }
        }
    }
    & & \\ \\

    \theutterance \stepcounter{utterance}  
    & & & \multicolumn{2}{p{0.3\linewidth}}{
        \cellcolor[rgb]{0.9,0.9,0.9}{
            \makecell[{{p{\linewidth}}}]{
                \texttt{\tiny{[GM$|$GM]}}
                \texttt{BEWEEG: T, (5,1)} \\
            }
        }
    }
    & & \\ \\

    \theutterance \stepcounter{utterance}  
    & & & \multicolumn{2}{p{0.3\linewidth}}{
        \cellcolor[rgb]{0.9,0.9,0.9}{
            \makecell[{{p{\linewidth}}}]{
                \texttt{\tiny{[GM$|$GM]}}
                \texttt{Straf: (5,1) is niet leeg, maar bevat '│'.} \\
\\ 
\texttt{Jullie hebben gezamenlijk 5 van 16 straffen verzameld.} \\
\texttt{Probeer het opnieuw!} \\
            }
        }
    }
    & & \\ \\

    \theutterance \stepcounter{utterance}  
    & & \multicolumn{4}{p{0.6\linewidth}}{
        \cellcolor[rgb]{0.9,0.9,0.9}{
            \makecell[{{p{\linewidth}}}]{
                \texttt{\tiny{[P1$\langle$GM]}}
                \texttt{Straf: (5,1) is niet leeg, maar bevat '│'.} \\
\\ 
\texttt{Jullie hebben gezamenlijk 5 van 16 straffen verzameld.} \\
\texttt{Probeer het opnieuw!} \\
            }
        }
    }
    & \\ \\

    \theutterance \stepcounter{utterance}  
    & \multicolumn{4}{p{0.6\linewidth}}{
        \cellcolor[rgb]{0.9,0.9,0.9}{
            \makecell[{{p{\linewidth}}}]{
                \texttt{\tiny{[P1$\rangle$GM]}}
                \texttt{BEWEEG: T, (1,1)} \\
            }
        }
    }
    & & \\ \\

    \theutterance \stepcounter{utterance}  
    & & & \multicolumn{2}{p{0.3\linewidth}}{
        \cellcolor[rgb]{0.9,0.9,0.9}{
            \makecell[{{p{\linewidth}}}]{
                \texttt{\tiny{[GM$|$GM]}}
                \texttt{BEWEEG: T, (1,1)} \\
            }
        }
    }
    & & \\ \\

    \theutterance \stepcounter{utterance}  
    & & & \multicolumn{2}{p{0.3\linewidth}}{
        \cellcolor[rgb]{0.9,0.9,0.9}{
            \makecell[{{p{\linewidth}}}]{
                \texttt{\tiny{[GM$|$GM]}}
                \texttt{T moved to (1, 1)} \\
            }
        }
    }
    & & \\ \\

    \theutterance \stepcounter{utterance}  
    & & \multicolumn{4}{p{0.6\linewidth}}{
        \cellcolor[rgb]{0.9,0.9,0.9}{
            \makecell[{{p{\linewidth}}}]{
                \texttt{\tiny{[GM$\rangle$P2]}}
                \texttt{Verplaatst 'T' naar (7,1) met succes. Uw bijgewerkte raster ziet er als volgt uit:} \\
\texttt{```} \\
\texttt{ 1234567} \\
\texttt{╔═╤══╤══╗} \\
\texttt{║◌│◌O│◌T║ 1} \\
\texttt{╟─┤◌◌│◌◌║ 2} \\
\texttt{║◌│◌◌├──╢ 3} \\
\texttt{╟─┘◌D│◌◌║ 4} \\
\texttt{║◌SA◌│◌◌║ 5} \\
\texttt{║PM┌─┼──╢ 6} \\
\texttt{║◌◌│◌│◌◌║ 7} \\
\texttt{╚══╧═╧══╝} \\
\texttt{```} \\
\texttt{Je speelt momenteel ronde 3 van maximaal 28.} \\
\texttt{Jullie hebben gezamenlijk 5 van 16 straffen verzameld.} \\
\texttt{De andere speler heeft een voorwerp op zijn rooster verplaatst.} \\
\\ 
\\ 
\texttt{Wat is je volgende opdracht?} \\
            }
        }
    }
    & \\ \\

    \theutterance \stepcounter{utterance}  
    & & & \multicolumn{4}{p{0.6\linewidth}}{
        \cellcolor[rgb]{0.9,0.9,0.9}{
            \makecell[{{p{\linewidth}}}]{
                \texttt{\tiny{[GM$\langle$P2]}}
                \texttt{ZEG: Goed gedaan met de T. Ik heb de T nu naar (7,1) verplaatst. Ik ga nu proberen de O naar de doelpositie (4,5) te verplaatsen.} \\
            }
        }
    }
     \\ \\

    \theutterance \stepcounter{utterance}  
    & & & \multicolumn{2}{p{0.3\linewidth}}{
        \cellcolor[rgb]{0.9,0.9,0.9}{
            \makecell[{{p{\linewidth}}}]{
                \texttt{\tiny{[GM$|$GM]}}
                \texttt{ZEG: Goed gedaan met de T. Ik heb de T nu naar (7,1) verplaatst. Ik ga nu proberen de O naar de doelpositie (4,5) te verplaatsen.} \\
            }
        }
    }
    & & \\ \\

    \theutterance \stepcounter{utterance}  
    & & \multicolumn{4}{p{0.6\linewidth}}{
        \cellcolor[rgb]{0.9,0.9,0.9}{
            \makecell[{{p{\linewidth}}}]{
                \texttt{\tiny{[P1$\langle$GM]}}
                \texttt{Verplaatst 'T' naar (1,1) met succes. Uw bijgewerkte raster ziet er als volgt uit:} \\
\texttt{```} \\
\texttt{ 1234567} \\
\texttt{╔═╤══╤══╗} \\
\texttt{║T│◌◌│◌◌║ 1} \\
\texttt{╟─┤◌◌│◌S║ 2} \\
\texttt{║◌│◌D├──╢ 3} \\
\texttt{╟─┘◌◌│O◌║ 4} \\
\texttt{║◌◌PM│◌A║ 5} \\
\texttt{║◌◌┌─┼──╢ 6} \\
\texttt{║◌◌│◌│◌◌║ 7} \\
\texttt{╚══╧═╧══╝} \\
\texttt{```} \\
\texttt{Je speelt momenteel ronde 3 van maximaal 28.} \\
\texttt{Jullie hebben gezamenlijk 5 van 16 straffen verzameld.} \\
\texttt{De andere speler stuurde dit bericht:} \\
\\ 
\\ 
\texttt{"Goed gedaan met de T. Ik heb de T nu naar (7,1) verplaatst. Ik ga nu proberen de O naar de doelpositie (4,5) te verplaatsen."} \\
\\ 
\\ 
\texttt{Wat is je volgende opdracht?} \\
            }
        }
    }
    & \\ \\

    \theutterance \stepcounter{utterance}  
    & \multicolumn{4}{p{0.6\linewidth}}{
        \cellcolor[rgb]{0.9,0.9,0.9}{
            \makecell[{{p{\linewidth}}}]{
                \texttt{\tiny{[P1$\rangle$GM]}}
                \texttt{BEWEEG: S, (4,4)} \\
            }
        }
    }
    & & \\ \\

    \theutterance \stepcounter{utterance}  
    & & & \multicolumn{2}{p{0.3\linewidth}}{
        \cellcolor[rgb]{0.9,0.9,0.9}{
            \makecell[{{p{\linewidth}}}]{
                \texttt{\tiny{[GM$|$GM]}}
                \texttt{BEWEEG: S, (4,4)} \\
            }
        }
    }
    & & \\ \\

    \theutterance \stepcounter{utterance}  
    & & & \multicolumn{2}{p{0.3\linewidth}}{
        \cellcolor[rgb]{0.9,0.9,0.9}{
            \makecell[{{p{\linewidth}}}]{
                \texttt{\tiny{[GM$|$GM]}}
                \texttt{S moved to (4, 4)} \\
            }
        }
    }
    & & \\ \\

    \theutterance \stepcounter{utterance}  
    & & \multicolumn{4}{p{0.6\linewidth}}{
        \cellcolor[rgb]{0.9,0.9,0.9}{
            \makecell[{{p{\linewidth}}}]{
                \texttt{\tiny{[GM$\rangle$P2]}}
                \texttt{Je bericht is doorgestuurd naar de andere speler.} \\
\\ 
\\ 
\texttt{Je speelt momenteel ronde 4 van maximaal 28.} \\
\texttt{Jullie hebben gezamenlijk 5 van 16 straffen verzameld.} \\
\texttt{De andere speler heeft een voorwerp op zijn rooster verplaatst.} \\
\\ 
\\ 
\texttt{Wat is je volgende opdracht?} \\
            }
        }
    }
    & \\ \\

    \theutterance \stepcounter{utterance}  
    & & & \multicolumn{4}{p{0.6\linewidth}}{
        \cellcolor[rgb]{0.9,0.9,0.9}{
            \makecell[{{p{\linewidth}}}]{
                \texttt{\tiny{[GM$\langle$P2]}}
                \texttt{BEWEEG: O, (4,5)} \\
            }
        }
    }
     \\ \\

    \theutterance \stepcounter{utterance}  
    & & & \multicolumn{2}{p{0.3\linewidth}}{
        \cellcolor[rgb]{0.9,0.9,0.9}{
            \makecell[{{p{\linewidth}}}]{
                \texttt{\tiny{[GM$|$GM]}}
                \texttt{BEWEEG: O, (4,5)} \\
            }
        }
    }
    & & \\ \\

    \theutterance \stepcounter{utterance}  
    & & & \multicolumn{2}{p{0.3\linewidth}}{
        \cellcolor[rgb]{0.9,0.9,0.9}{
            \makecell[{{p{\linewidth}}}]{
                \texttt{\tiny{[GM$|$GM]}}
                \texttt{O moved to (4, 5)} \\
            }
        }
    }
    & & \\ \\

    \theutterance \stepcounter{utterance}  
    & & \multicolumn{4}{p{0.6\linewidth}}{
        \cellcolor[rgb]{0.9,0.9,0.9}{
            \makecell[{{p{\linewidth}}}]{
                \texttt{\tiny{[P1$\langle$GM]}}
                \texttt{Verplaatst 'S' naar (4,4) met succes. Uw bijgewerkte raster ziet er als volgt uit:} \\
\texttt{```} \\
\texttt{ 1234567} \\
\texttt{╔═╤══╤══╗} \\
\texttt{║T│◌◌│◌◌║ 1} \\
\texttt{╟─┤◌◌│◌◌║ 2} \\
\texttt{║◌│◌D├──╢ 3} \\
\texttt{╟─┘◌S│O◌║ 4} \\
\texttt{║◌◌PM│◌A║ 5} \\
\texttt{║◌◌┌─┼──╢ 6} \\
\texttt{║◌◌│◌│◌◌║ 7} \\
\texttt{╚══╧═╧══╝} \\
\texttt{```} \\
\texttt{Je speelt momenteel ronde 4 van maximaal 28.} \\
\texttt{Jullie hebben gezamenlijk 5 van 16 straffen verzameld.} \\
\texttt{De andere speler heeft een voorwerp op zijn rooster verplaatst.} \\
\\ 
\\ 
\texttt{Wat is je volgende opdracht?} \\
            }
        }
    }
    & \\ \\

    \theutterance \stepcounter{utterance}  
    & \multicolumn{4}{p{0.6\linewidth}}{
        \cellcolor[rgb]{0.9,0.9,0.9}{
            \makecell[{{p{\linewidth}}}]{
                \texttt{\tiny{[P1$\rangle$GM]}}
                \texttt{ZEG: Goed bezig! Ik heb S nu naar (4,4) verplaatst. Ik ga nu proberen om mijn O naar de doelpositie (4,5) te verplaatsen.} \\
            }
        }
    }
    & & \\ \\

    \theutterance \stepcounter{utterance}  
    & & & \multicolumn{2}{p{0.3\linewidth}}{
        \cellcolor[rgb]{0.9,0.9,0.9}{
            \makecell[{{p{\linewidth}}}]{
                \texttt{\tiny{[GM$|$GM]}}
                \texttt{ZEG: Goed bezig! Ik heb S nu naar (4,4) verplaatst. Ik ga nu proberen om mijn O naar de doelpositie (4,5) te verplaatsen.} \\
            }
        }
    }
    & & \\ \\

    \theutterance \stepcounter{utterance}  
    & & \multicolumn{4}{p{0.6\linewidth}}{
        \cellcolor[rgb]{0.9,0.9,0.9}{
            \makecell[{{p{\linewidth}}}]{
                \texttt{\tiny{[GM$\rangle$P2]}}
                \texttt{Verplaatst 'O' naar (4,5) met succes. Uw bijgewerkte raster ziet er als volgt uit:} \\
\texttt{```} \\
\texttt{ 1234567} \\
\texttt{╔═╤══╤══╗} \\
\texttt{║◌│◌◌│◌T║ 1} \\
\texttt{╟─┤◌◌│◌◌║ 2} \\
\texttt{║◌│◌◌├──╢ 3} \\
\texttt{╟─┘◌D│◌◌║ 4} \\
\texttt{║◌SAO│◌◌║ 5} \\
\texttt{║PM┌─┼──╢ 6} \\
\texttt{║◌◌│◌│◌◌║ 7} \\
\texttt{╚══╧═╧══╝} \\
\texttt{```} \\
\texttt{Je speelt momenteel ronde 5 van maximaal 28.} \\
\texttt{Jullie hebben gezamenlijk 5 van 16 straffen verzameld.} \\
\texttt{De andere speler stuurde dit bericht:} \\
\\ 
\\ 
\texttt{"Goed bezig! Ik heb S nu naar (4,4) verplaatst. Ik ga nu proberen om mijn O naar de doelpositie (4,5) te verplaatsen."} \\
\\ 
\\ 
\texttt{Wat is je volgende opdracht?} \\
            }
        }
    }
    & \\ \\

    \theutterance \stepcounter{utterance}  
    & & & \multicolumn{4}{p{0.6\linewidth}}{
        \cellcolor[rgb]{0.9,0.9,0.9}{
            \makecell[{{p{\linewidth}}}]{
                \texttt{\tiny{[GM$\langle$P2]}}
                \texttt{ZEG: Dat is goed! Ik heb mijn O al naar (4,5) verplaatst. Ik ga nu mijn D naar de doelpositie (3,5) verplaatsen.} \\
            }
        }
    }
     \\ \\

    \theutterance \stepcounter{utterance}  
    & & & \multicolumn{2}{p{0.3\linewidth}}{
        \cellcolor[rgb]{0.9,0.9,0.9}{
            \makecell[{{p{\linewidth}}}]{
                \texttt{\tiny{[GM$|$GM]}}
                \texttt{ZEG: Dat is goed! Ik heb mijn O al naar (4,5) verplaatst. Ik ga nu mijn D naar de doelpositie (3,5) verplaatsen.} \\
            }
        }
    }
    & & \\ \\

    \theutterance \stepcounter{utterance}  
    & & \multicolumn{4}{p{0.6\linewidth}}{
        \cellcolor[rgb]{0.9,0.9,0.9}{
            \makecell[{{p{\linewidth}}}]{
                \texttt{\tiny{[P1$\langle$GM]}}
                \texttt{Je bericht is doorgestuurd naar de andere speler.} \\
\\ 
\\ 
\texttt{Je speelt momenteel ronde 5 van maximaal 28.} \\
\texttt{Jullie hebben gezamenlijk 5 van 16 straffen verzameld.} \\
\texttt{De andere speler stuurde dit bericht:} \\
\\ 
\\ 
\texttt{"Dat is goed! Ik heb mijn O al naar (4,5) verplaatst. Ik ga nu mijn D naar de doelpositie (3,5) verplaatsen."} \\
\\ 
\\ 
\texttt{Wat is je volgende opdracht?} \\
            }
        }
    }
    & \\ \\

    \theutterance \stepcounter{utterance}  
    & \multicolumn{4}{p{0.6\linewidth}}{
        \cellcolor[rgb]{0.9,0.9,0.9}{
            \makecell[{{p{\linewidth}}}]{
                \texttt{\tiny{[P1$\rangle$GM]}}
                \texttt{BEWEEG: O, (4,5)} \\
            }
        }
    }
    & & \\ \\

    \theutterance \stepcounter{utterance}  
    & & & \multicolumn{2}{p{0.3\linewidth}}{
        \cellcolor[rgb]{0.9,0.9,0.9}{
            \makecell[{{p{\linewidth}}}]{
                \texttt{\tiny{[GM$|$GM]}}
                \texttt{BEWEEG: O, (4,5)} \\
            }
        }
    }
    & & \\ \\

    \theutterance \stepcounter{utterance}  
    & & & \multicolumn{2}{p{0.3\linewidth}}{
        \cellcolor[rgb]{0.9,0.9,0.9}{
            \makecell[{{p{\linewidth}}}]{
                \texttt{\tiny{[GM$|$GM]}}
                \texttt{Straf: (4,5) is niet leeg, maar bevat 'M'.} \\
\\ 
\texttt{Jullie hebben gezamenlijk 6 van 16 straffen verzameld.} \\
\texttt{Probeer het opnieuw!} \\
            }
        }
    }
    & & \\ \\

    \theutterance \stepcounter{utterance}  
    & & \multicolumn{4}{p{0.6\linewidth}}{
        \cellcolor[rgb]{0.9,0.9,0.9}{
            \makecell[{{p{\linewidth}}}]{
                \texttt{\tiny{[P1$\langle$GM]}}
                \texttt{Straf: (4,5) is niet leeg, maar bevat 'M'.} \\
\\ 
\texttt{Jullie hebben gezamenlijk 6 van 16 straffen verzameld.} \\
\texttt{Probeer het opnieuw!} \\
            }
        }
    }
    & \\ \\

    \theutterance \stepcounter{utterance}  
    & \multicolumn{4}{p{0.6\linewidth}}{
        \cellcolor[rgb]{0.9,0.9,0.9}{
            \makecell[{{p{\linewidth}}}]{
                \texttt{\tiny{[P1$\rangle$GM]}}
                \texttt{BEWEEG: O, (6,5)} \\
            }
        }
    }
    & & \\ \\

    \theutterance \stepcounter{utterance}  
    & & & \multicolumn{2}{p{0.3\linewidth}}{
        \cellcolor[rgb]{0.9,0.9,0.9}{
            \makecell[{{p{\linewidth}}}]{
                \texttt{\tiny{[GM$|$GM]}}
                \texttt{BEWEEG: O, (6,5)} \\
            }
        }
    }
    & & \\ \\

    \theutterance \stepcounter{utterance}  
    & & & \multicolumn{2}{p{0.3\linewidth}}{
        \cellcolor[rgb]{0.9,0.9,0.9}{
            \makecell[{{p{\linewidth}}}]{
                \texttt{\tiny{[GM$|$GM]}}
                \texttt{O moved to (6, 5)} \\
            }
        }
    }
    & & \\ \\

    \theutterance \stepcounter{utterance}  
    & & \multicolumn{4}{p{0.6\linewidth}}{
        \cellcolor[rgb]{0.9,0.9,0.9}{
            \makecell[{{p{\linewidth}}}]{
                \texttt{\tiny{[GM$\rangle$P2]}}
                \texttt{Je bericht is doorgestuurd naar de andere speler.} \\
\\ 
\\ 
\texttt{Je speelt momenteel ronde 6 van maximaal 28.} \\
\texttt{Jullie hebben gezamenlijk 6 van 16 straffen verzameld.} \\
\texttt{De andere speler heeft een voorwerp op zijn rooster verplaatst.} \\
\\ 
\\ 
\texttt{Wat is je volgende opdracht?} \\
            }
        }
    }
    & \\ \\

    \theutterance \stepcounter{utterance}  
    & & & \multicolumn{4}{p{0.6\linewidth}}{
        \cellcolor[rgb]{0.9,0.9,0.9}{
            \makecell[{{p{\linewidth}}}]{
                \texttt{\tiny{[GM$\langle$P2]}}
                \texttt{BEWEEG: D, (3,5)} \\
            }
        }
    }
     \\ \\

    \theutterance \stepcounter{utterance}  
    & & & \multicolumn{2}{p{0.3\linewidth}}{
        \cellcolor[rgb]{0.9,0.9,0.9}{
            \makecell[{{p{\linewidth}}}]{
                \texttt{\tiny{[GM$|$GM]}}
                \texttt{BEWEEG: D, (3,5)} \\
            }
        }
    }
    & & \\ \\

    \theutterance \stepcounter{utterance}  
    & & & \multicolumn{2}{p{0.3\linewidth}}{
        \cellcolor[rgb]{0.9,0.9,0.9}{
            \makecell[{{p{\linewidth}}}]{
                \texttt{\tiny{[GM$|$GM]}}
                \texttt{Straf: (3,5) is niet leeg, maar bevat 'A'.} \\
\\ 
\texttt{Jullie hebben gezamenlijk 7 van 16 straffen verzameld.} \\
\texttt{Probeer het opnieuw!} \\
            }
        }
    }
    & & \\ \\

    \theutterance \stepcounter{utterance}  
    & & \multicolumn{4}{p{0.6\linewidth}}{
        \cellcolor[rgb]{0.9,0.9,0.9}{
            \makecell[{{p{\linewidth}}}]{
                \texttt{\tiny{[GM$\rangle$P2]}}
                \texttt{Straf: (3,5) is niet leeg, maar bevat 'A'.} \\
\\ 
\texttt{Jullie hebben gezamenlijk 7 van 16 straffen verzameld.} \\
\texttt{Probeer het opnieuw!} \\
            }
        }
    }
    & \\ \\

    \theutterance \stepcounter{utterance}  
    & & & \multicolumn{4}{p{0.6\linewidth}}{
        \cellcolor[rgb]{0.9,0.9,0.9}{
            \makecell[{{p{\linewidth}}}]{
                \texttt{\tiny{[GM$\langle$P2]}}
                \texttt{BEWEEG: D, (4,3)} \\
            }
        }
    }
     \\ \\

    \theutterance \stepcounter{utterance}  
    & & & \multicolumn{2}{p{0.3\linewidth}}{
        \cellcolor[rgb]{0.9,0.9,0.9}{
            \makecell[{{p{\linewidth}}}]{
                \texttt{\tiny{[GM$|$GM]}}
                \texttt{BEWEEG: D, (4,3)} \\
            }
        }
    }
    & & \\ \\

    \theutterance \stepcounter{utterance}  
    & & & \multicolumn{2}{p{0.3\linewidth}}{
        \cellcolor[rgb]{0.9,0.9,0.9}{
            \makecell[{{p{\linewidth}}}]{
                \texttt{\tiny{[GM$|$GM]}}
                \texttt{D moved to (4, 3)} \\
            }
        }
    }
    & & \\ \\

    \theutterance \stepcounter{utterance}  
    & & \multicolumn{4}{p{0.6\linewidth}}{
        \cellcolor[rgb]{0.9,0.9,0.9}{
            \makecell[{{p{\linewidth}}}]{
                \texttt{\tiny{[P1$\langle$GM]}}
                \texttt{Verplaatst 'O' naar (6,5) met succes. Uw bijgewerkte raster ziet er als volgt uit:} \\
\texttt{```} \\
\texttt{ 1234567} \\
\texttt{╔═╤══╤══╗} \\
\texttt{║T│◌◌│◌◌║ 1} \\
\texttt{╟─┤◌◌│◌◌║ 2} \\
\texttt{║◌│◌D├──╢ 3} \\
\texttt{╟─┘◌S│◌◌║ 4} \\
\texttt{║◌◌PM│OA║ 5} \\
\texttt{║◌◌┌─┼──╢ 6} \\
\texttt{║◌◌│◌│◌◌║ 7} \\
\texttt{╚══╧═╧══╝} \\
\texttt{```} \\
\texttt{Je speelt momenteel ronde 6 van maximaal 28.} \\
\texttt{Jullie hebben gezamenlijk 7 van 16 straffen verzameld.} \\
\texttt{De andere speler heeft een voorwerp op zijn rooster verplaatst.} \\
\\ 
\\ 
\texttt{Wat is je volgende opdracht?} \\
            }
        }
    }
    & \\ \\

    \theutterance \stepcounter{utterance}  
    & \multicolumn{4}{p{0.6\linewidth}}{
        \cellcolor[rgb]{0.9,0.9,0.9}{
            \makecell[{{p{\linewidth}}}]{
                \texttt{\tiny{[P1$\rangle$GM]}}
                \texttt{BEWEEG: O, (4,5)} \\
            }
        }
    }
    & & \\ \\

    \theutterance \stepcounter{utterance}  
    & & & \multicolumn{2}{p{0.3\linewidth}}{
        \cellcolor[rgb]{0.9,0.9,0.9}{
            \makecell[{{p{\linewidth}}}]{
                \texttt{\tiny{[GM$|$GM]}}
                \texttt{BEWEEG: O, (4,5)} \\
            }
        }
    }
    & & \\ \\

    \theutterance \stepcounter{utterance}  
    & & & \multicolumn{2}{p{0.3\linewidth}}{
        \cellcolor[rgb]{0.9,0.9,0.9}{
            \makecell[{{p{\linewidth}}}]{
                \texttt{\tiny{[GM$|$GM]}}
                \texttt{Straf: (4,5) is niet leeg, maar bevat 'M'.} \\
\\ 
\texttt{Jullie hebben gezamenlijk 8 van 16 straffen verzameld.} \\
\texttt{Probeer het opnieuw!} \\
            }
        }
    }
    & & \\ \\

    \theutterance \stepcounter{utterance}  
    & & \multicolumn{4}{p{0.6\linewidth}}{
        \cellcolor[rgb]{0.9,0.9,0.9}{
            \makecell[{{p{\linewidth}}}]{
                \texttt{\tiny{[P1$\langle$GM]}}
                \texttt{Straf: (4,5) is niet leeg, maar bevat 'M'.} \\
\\ 
\texttt{Jullie hebben gezamenlijk 8 van 16 straffen verzameld.} \\
\texttt{Probeer het opnieuw!} \\
            }
        }
    }
    & \\ \\

    \theutterance \stepcounter{utterance}  
    & \multicolumn{4}{p{0.6\linewidth}}{
        \cellcolor[rgb]{0.9,0.9,0.9}{
            \makecell[{{p{\linewidth}}}]{
                \texttt{\tiny{[P1$\rangle$GM]}}
                \texttt{BEWEEG: O, (5,4)} \\
            }
        }
    }
    & & \\ \\

    \theutterance \stepcounter{utterance}  
    & & & \multicolumn{2}{p{0.3\linewidth}}{
        \cellcolor[rgb]{0.9,0.9,0.9}{
            \makecell[{{p{\linewidth}}}]{
                \texttt{\tiny{[GM$|$GM]}}
                \texttt{BEWEEG: O, (5,4)} \\
            }
        }
    }
    & & \\ \\

    \theutterance \stepcounter{utterance}  
    & & & \multicolumn{2}{p{0.3\linewidth}}{
        \cellcolor[rgb]{0.9,0.9,0.9}{
            \makecell[{{p{\linewidth}}}]{
                \texttt{\tiny{[GM$|$GM]}}
                \texttt{Straf: (5,4) is niet leeg, maar bevat '│'.} \\
\\ 
\texttt{Jullie hebben gezamenlijk 9 van 16 straffen verzameld.} \\
\texttt{Probeer het opnieuw!} \\
            }
        }
    }
    & & \\ \\

    \theutterance \stepcounter{utterance}  
    & & \multicolumn{4}{p{0.6\linewidth}}{
        \cellcolor[rgb]{0.9,0.9,0.9}{
            \makecell[{{p{\linewidth}}}]{
                \texttt{\tiny{[P1$\langle$GM]}}
                \texttt{Straf: (5,4) is niet leeg, maar bevat '│'.} \\
\\ 
\texttt{Jullie hebben gezamenlijk 9 van 16 straffen verzameld.} \\
\texttt{Probeer het opnieuw!} \\
            }
        }
    }
    & \\ \\

    \theutterance \stepcounter{utterance}  
    & \multicolumn{4}{p{0.6\linewidth}}{
        \cellcolor[rgb]{0.9,0.9,0.9}{
            \makecell[{{p{\linewidth}}}]{
                \texttt{\tiny{[P1$\rangle$GM]}}
                \texttt{BEWEEG: O, (6,4)} \\
            }
        }
    }
    & & \\ \\

    \theutterance \stepcounter{utterance}  
    & & & \multicolumn{2}{p{0.3\linewidth}}{
        \cellcolor[rgb]{0.9,0.9,0.9}{
            \makecell[{{p{\linewidth}}}]{
                \texttt{\tiny{[GM$|$GM]}}
                \texttt{BEWEEG: O, (6,4)} \\
            }
        }
    }
    & & \\ \\

    \theutterance \stepcounter{utterance}  
    & & & \multicolumn{2}{p{0.3\linewidth}}{
        \cellcolor[rgb]{0.9,0.9,0.9}{
            \makecell[{{p{\linewidth}}}]{
                \texttt{\tiny{[GM$|$GM]}}
                \texttt{O moved to (6, 4)} \\
            }
        }
    }
    & & \\ \\

    \theutterance \stepcounter{utterance}  
    & & \multicolumn{4}{p{0.6\linewidth}}{
        \cellcolor[rgb]{0.9,0.9,0.9}{
            \makecell[{{p{\linewidth}}}]{
                \texttt{\tiny{[GM$\rangle$P2]}}
                \texttt{Verplaatst 'D' naar (4,3) met succes. Uw bijgewerkte raster ziet er als volgt uit:} \\
\texttt{```} \\
\texttt{ 1234567} \\
\texttt{╔═╤══╤══╗} \\
\texttt{║◌│◌◌│◌T║ 1} \\
\texttt{╟─┤◌◌│◌◌║ 2} \\
\texttt{║◌│◌D├──╢ 3} \\
\texttt{╟─┘◌◌│◌◌║ 4} \\
\texttt{║◌SAO│◌◌║ 5} \\
\texttt{║PM┌─┼──╢ 6} \\
\texttt{║◌◌│◌│◌◌║ 7} \\
\texttt{╚══╧═╧══╝} \\
\texttt{```} \\
\texttt{Je speelt momenteel ronde 7 van maximaal 28.} \\
\texttt{Jullie hebben gezamenlijk 9 van 16 straffen verzameld.} \\
\texttt{De andere speler heeft een voorwerp op zijn rooster verplaatst.} \\
\\ 
\\ 
\texttt{Wat is je volgende opdracht?} \\
            }
        }
    }
    & \\ \\

    \theutterance \stepcounter{utterance}  
    & & & \multicolumn{4}{p{0.6\linewidth}}{
        \cellcolor[rgb]{0.9,0.9,0.9}{
            \makecell[{{p{\linewidth}}}]{
                \texttt{\tiny{[GM$\langle$P2]}}
                \texttt{ZEG: Ik heb D nu naar (4,3) verplaatst. Ik ga nu proberen de S naar (3,4) te verplaatsen.} \\
            }
        }
    }
     \\ \\

    \theutterance \stepcounter{utterance}  
    & & & \multicolumn{2}{p{0.3\linewidth}}{
        \cellcolor[rgb]{0.9,0.9,0.9}{
            \makecell[{{p{\linewidth}}}]{
                \texttt{\tiny{[GM$|$GM]}}
                \texttt{ZEG: Ik heb D nu naar (4,3) verplaatst. Ik ga nu proberen de S naar (3,4) te verplaatsen.} \\
            }
        }
    }
    & & \\ \\

    \theutterance \stepcounter{utterance}  
    & & \multicolumn{4}{p{0.6\linewidth}}{
        \cellcolor[rgb]{0.9,0.9,0.9}{
            \makecell[{{p{\linewidth}}}]{
                \texttt{\tiny{[P1$\langle$GM]}}
                \texttt{Verplaatst 'O' naar (6,4) met succes. Uw bijgewerkte raster ziet er als volgt uit:} \\
\texttt{```} \\
\texttt{ 1234567} \\
\texttt{╔═╤══╤══╗} \\
\texttt{║T│◌◌│◌◌║ 1} \\
\texttt{╟─┤◌◌│◌◌║ 2} \\
\texttt{║◌│◌D├──╢ 3} \\
\texttt{╟─┘◌S│O◌║ 4} \\
\texttt{║◌◌PM│◌A║ 5} \\
\texttt{║◌◌┌─┼──╢ 6} \\
\texttt{║◌◌│◌│◌◌║ 7} \\
\texttt{╚══╧═╧══╝} \\
\texttt{```} \\
\texttt{Je speelt momenteel ronde 7 van maximaal 28.} \\
\texttt{Jullie hebben gezamenlijk 9 van 16 straffen verzameld.} \\
\texttt{De andere speler stuurde dit bericht:} \\
\\ 
\\ 
\texttt{"Ik heb D nu naar (4,3) verplaatst. Ik ga nu proberen de S naar (3,4) te verplaatsen."} \\
\\ 
\\ 
\texttt{Wat is je volgende opdracht?} \\
            }
        }
    }
    & \\ \\

    \theutterance \stepcounter{utterance}  
    & \multicolumn{4}{p{0.6\linewidth}}{
        \cellcolor[rgb]{0.9,0.9,0.9}{
            \makecell[{{p{\linewidth}}}]{
                \texttt{\tiny{[P1$\rangle$GM]}}
                \texttt{BEWEEG: P, (3,3)} \\
            }
        }
    }
    & & \\ \\

    \theutterance \stepcounter{utterance}  
    & & & \multicolumn{2}{p{0.3\linewidth}}{
        \cellcolor[rgb]{0.9,0.9,0.9}{
            \makecell[{{p{\linewidth}}}]{
                \texttt{\tiny{[GM$|$GM]}}
                \texttt{BEWEEG: P, (3,3)} \\
            }
        }
    }
    & & \\ \\

    \theutterance \stepcounter{utterance}  
    & & & \multicolumn{2}{p{0.3\linewidth}}{
        \cellcolor[rgb]{0.9,0.9,0.9}{
            \makecell[{{p{\linewidth}}}]{
                \texttt{\tiny{[GM$|$GM]}}
                \texttt{P moved to (3, 3)} \\
            }
        }
    }
    & & \\ \\

    \theutterance \stepcounter{utterance}  
    & & \multicolumn{4}{p{0.6\linewidth}}{
        \cellcolor[rgb]{0.9,0.9,0.9}{
            \makecell[{{p{\linewidth}}}]{
                \texttt{\tiny{[GM$\rangle$P2]}}
                \texttt{Je bericht is doorgestuurd naar de andere speler.} \\
\\ 
\\ 
\texttt{Je speelt momenteel ronde 8 van maximaal 28.} \\
\texttt{Jullie hebben gezamenlijk 9 van 16 straffen verzameld.} \\
\texttt{De andere speler heeft een voorwerp op zijn rooster verplaatst.} \\
\\ 
\\ 
\texttt{Wat is je volgende opdracht?} \\
            }
        }
    }
    & \\ \\

    \theutterance \stepcounter{utterance}  
    & & & \multicolumn{4}{p{0.6\linewidth}}{
        \cellcolor[rgb]{0.9,0.9,0.9}{
            \makecell[{{p{\linewidth}}}]{
                \texttt{\tiny{[GM$\langle$P2]}}
                \texttt{BEWEEG: S, (3,4)} \\
            }
        }
    }
     \\ \\

    \theutterance \stepcounter{utterance}  
    & & & \multicolumn{2}{p{0.3\linewidth}}{
        \cellcolor[rgb]{0.9,0.9,0.9}{
            \makecell[{{p{\linewidth}}}]{
                \texttt{\tiny{[GM$|$GM]}}
                \texttt{BEWEEG: S, (3,4)} \\
            }
        }
    }
    & & \\ \\

    \theutterance \stepcounter{utterance}  
    & & & \multicolumn{2}{p{0.3\linewidth}}{
        \cellcolor[rgb]{0.9,0.9,0.9}{
            \makecell[{{p{\linewidth}}}]{
                \texttt{\tiny{[GM$|$GM]}}
                \texttt{S moved to (3, 4)} \\
            }
        }
    }
    & & \\ \\

    \theutterance \stepcounter{utterance}  
    & & \multicolumn{4}{p{0.6\linewidth}}{
        \cellcolor[rgb]{0.9,0.9,0.9}{
            \makecell[{{p{\linewidth}}}]{
                \texttt{\tiny{[P1$\langle$GM]}}
                \texttt{Verplaatst 'P' naar (3,3) met succes. Uw bijgewerkte raster ziet er als volgt uit:} \\
\texttt{```} \\
\texttt{ 1234567} \\
\texttt{╔═╤══╤══╗} \\
\texttt{║T│◌◌│◌◌║ 1} \\
\texttt{╟─┤◌◌│◌◌║ 2} \\
\texttt{║◌│PD├──╢ 3} \\
\texttt{╟─┘◌S│O◌║ 4} \\
\texttt{║◌◌◌M│◌A║ 5} \\
\texttt{║◌◌┌─┼──╢ 6} \\
\texttt{║◌◌│◌│◌◌║ 7} \\
\texttt{╚══╧═╧══╝} \\
\texttt{```} \\
\texttt{Je speelt momenteel ronde 8 van maximaal 28.} \\
\texttt{Jullie hebben gezamenlijk 9 van 16 straffen verzameld.} \\
\texttt{De andere speler heeft een voorwerp op zijn rooster verplaatst.} \\
\\ 
\\ 
\texttt{Wat is je volgende opdracht?} \\
            }
        }
    }
    & \\ \\

    \theutterance \stepcounter{utterance}  
    & \multicolumn{4}{p{0.6\linewidth}}{
        \cellcolor[rgb]{0.9,0.9,0.9}{
            \makecell[{{p{\linewidth}}}]{
                \texttt{\tiny{[P1$\rangle$GM]}}
                \texttt{BEWEEG: M, (4,3)} \\
            }
        }
    }
    & & \\ \\

    \theutterance \stepcounter{utterance}  
    & & & \multicolumn{2}{p{0.3\linewidth}}{
        \cellcolor[rgb]{0.9,0.9,0.9}{
            \makecell[{{p{\linewidth}}}]{
                \texttt{\tiny{[GM$|$GM]}}
                \texttt{BEWEEG: M, (4,3)} \\
            }
        }
    }
    & & \\ \\

    \theutterance \stepcounter{utterance}  
    & & & \multicolumn{2}{p{0.3\linewidth}}{
        \cellcolor[rgb]{0.9,0.9,0.9}{
            \makecell[{{p{\linewidth}}}]{
                \texttt{\tiny{[GM$|$GM]}}
                \texttt{Straf: (4,3) is niet leeg, maar bevat 'D'.} \\
\\ 
\texttt{Jullie hebben gezamenlijk 10 van 16 straffen verzameld.} \\
\texttt{Probeer het opnieuw!} \\
            }
        }
    }
    & & \\ \\

    \theutterance \stepcounter{utterance}  
    & & \multicolumn{4}{p{0.6\linewidth}}{
        \cellcolor[rgb]{0.9,0.9,0.9}{
            \makecell[{{p{\linewidth}}}]{
                \texttt{\tiny{[P1$\langle$GM]}}
                \texttt{Straf: (4,3) is niet leeg, maar bevat 'D'.} \\
\\ 
\texttt{Jullie hebben gezamenlijk 10 van 16 straffen verzameld.} \\
\texttt{Probeer het opnieuw!} \\
            }
        }
    }
    & \\ \\

    \theutterance \stepcounter{utterance}  
    & \multicolumn{4}{p{0.6\linewidth}}{
        \cellcolor[rgb]{0.9,0.9,0.9}{
            \makecell[{{p{\linewidth}}}]{
                \texttt{\tiny{[P1$\rangle$GM]}}
                \texttt{BEWEEG: M, (4,5)} \\
            }
        }
    }
    & & \\ \\

    \theutterance \stepcounter{utterance}  
    & & & \multicolumn{2}{p{0.3\linewidth}}{
        \cellcolor[rgb]{0.9,0.9,0.9}{
            \makecell[{{p{\linewidth}}}]{
                \texttt{\tiny{[GM$|$GM]}}
                \texttt{BEWEEG: M, (4,5)} \\
            }
        }
    }
    & & \\ \\

    \theutterance \stepcounter{utterance}  
    & & & \multicolumn{2}{p{0.3\linewidth}}{
        \cellcolor[rgb]{0.9,0.9,0.9}{
            \makecell[{{p{\linewidth}}}]{
                \texttt{\tiny{[GM$|$GM]}}
                \texttt{Straf: (4,5) is niet leeg, maar bevat 'M'.} \\
\\ 
\texttt{Jullie hebben gezamenlijk 11 van 16 straffen verzameld.} \\
\texttt{Probeer het opnieuw!} \\
            }
        }
    }
    & & \\ \\

    \theutterance \stepcounter{utterance}  
    & & \multicolumn{4}{p{0.6\linewidth}}{
        \cellcolor[rgb]{0.9,0.9,0.9}{
            \makecell[{{p{\linewidth}}}]{
                \texttt{\tiny{[P1$\langle$GM]}}
                \texttt{Straf: (4,5) is niet leeg, maar bevat 'M'.} \\
\\ 
\texttt{Jullie hebben gezamenlijk 11 van 16 straffen verzameld.} \\
\texttt{Probeer het opnieuw!} \\
            }
        }
    }
    & \\ \\

    \theutterance \stepcounter{utterance}  
    & \multicolumn{4}{p{0.6\linewidth}}{
        \cellcolor[rgb]{0.9,0.9,0.9}{
            \makecell[{{p{\linewidth}}}]{
                \texttt{\tiny{[P1$\rangle$GM]}}
                \texttt{BEWEEG: M, (4,4)} \\
            }
        }
    }
    & & \\ \\

    \theutterance \stepcounter{utterance}  
    & & & \multicolumn{2}{p{0.3\linewidth}}{
        \cellcolor[rgb]{0.9,0.9,0.9}{
            \makecell[{{p{\linewidth}}}]{
                \texttt{\tiny{[GM$|$GM]}}
                \texttt{BEWEEG: M, (4,4)} \\
            }
        }
    }
    & & \\ \\

    \theutterance \stepcounter{utterance}  
    & & & \multicolumn{2}{p{0.3\linewidth}}{
        \cellcolor[rgb]{0.9,0.9,0.9}{
            \makecell[{{p{\linewidth}}}]{
                \texttt{\tiny{[GM$|$GM]}}
                \texttt{Straf: (4,4) is niet leeg, maar bevat 'S'.} \\
\\ 
\texttt{Jullie hebben gezamenlijk 12 van 16 straffen verzameld.} \\
\texttt{Probeer het opnieuw!} \\
            }
        }
    }
    & & \\ \\

    \theutterance \stepcounter{utterance}  
    & & \multicolumn{4}{p{0.6\linewidth}}{
        \cellcolor[rgb]{0.9,0.9,0.9}{
            \makecell[{{p{\linewidth}}}]{
                \texttt{\tiny{[P1$\langle$GM]}}
                \texttt{Straf: (4,4) is niet leeg, maar bevat 'S'.} \\
\\ 
\texttt{Jullie hebben gezamenlijk 12 van 16 straffen verzameld.} \\
\texttt{Probeer het opnieuw!} \\
            }
        }
    }
    & \\ \\

    \theutterance \stepcounter{utterance}  
    & \multicolumn{4}{p{0.6\linewidth}}{
        \cellcolor[rgb]{0.9,0.9,0.9}{
            \makecell[{{p{\linewidth}}}]{
                \texttt{\tiny{[P1$\rangle$GM]}}
                \texttt{BEWEEG: M, (4,2)} \\
            }
        }
    }
    & & \\ \\

    \theutterance \stepcounter{utterance}  
    & & & \multicolumn{2}{p{0.3\linewidth}}{
        \cellcolor[rgb]{0.9,0.9,0.9}{
            \makecell[{{p{\linewidth}}}]{
                \texttt{\tiny{[GM$|$GM]}}
                \texttt{BEWEEG: M, (4,2)} \\
            }
        }
    }
    & & \\ \\

    \theutterance \stepcounter{utterance}  
    & & & \multicolumn{2}{p{0.3\linewidth}}{
        \cellcolor[rgb]{0.9,0.9,0.9}{
            \makecell[{{p{\linewidth}}}]{
                \texttt{\tiny{[GM$|$GM]}}
                \texttt{M moved to (4, 2)} \\
            }
        }
    }
    & & \\ \\

    \theutterance \stepcounter{utterance}  
    & & \multicolumn{4}{p{0.6\linewidth}}{
        \cellcolor[rgb]{0.9,0.9,0.9}{
            \makecell[{{p{\linewidth}}}]{
                \texttt{\tiny{[GM$\rangle$P2]}}
                \texttt{Verplaatst 'S' naar (3,4) met succes. Uw bijgewerkte raster ziet er als volgt uit:} \\
\texttt{```} \\
\texttt{ 1234567} \\
\texttt{╔═╤══╤══╗} \\
\texttt{║◌│◌◌│◌T║ 1} \\
\texttt{╟─┤◌◌│◌◌║ 2} \\
\texttt{║◌│◌D├──╢ 3} \\
\texttt{╟─┘S◌│◌◌║ 4} \\
\texttt{║◌◌AO│◌◌║ 5} \\
\texttt{║PM┌─┼──╢ 6} \\
\texttt{║◌◌│◌│◌◌║ 7} \\
\texttt{╚══╧═╧══╝} \\
\texttt{```} \\
\texttt{Je speelt momenteel ronde 9 van maximaal 28.} \\
\texttt{Jullie hebben gezamenlijk 12 van 16 straffen verzameld.} \\
\texttt{De andere speler heeft een voorwerp op zijn rooster verplaatst.} \\
\\ 
\\ 
\texttt{Wat is je volgende opdracht?} \\
            }
        }
    }
    & \\ \\

    \theutterance \stepcounter{utterance}  
    & & & \multicolumn{4}{p{0.6\linewidth}}{
        \cellcolor[rgb]{0.9,0.9,0.9}{
            \makecell[{{p{\linewidth}}}]{
                \texttt{\tiny{[GM$\langle$P2]}}
                \texttt{ZEG: Ik heb S nu naar (3,4) verplaatst. Ik ga nu proberen de A naar de doelpositie (4,4) te verplaatsen.} \\
            }
        }
    }
     \\ \\

    \theutterance \stepcounter{utterance}  
    & & & \multicolumn{2}{p{0.3\linewidth}}{
        \cellcolor[rgb]{0.9,0.9,0.9}{
            \makecell[{{p{\linewidth}}}]{
                \texttt{\tiny{[GM$|$GM]}}
                \texttt{ZEG: Ik heb S nu naar (3,4) verplaatst. Ik ga nu proberen de A naar de doelpositie (4,4) te verplaatsen.} \\
            }
        }
    }
    & & \\ \\

    \theutterance \stepcounter{utterance}  
    & & \multicolumn{4}{p{0.6\linewidth}}{
        \cellcolor[rgb]{0.9,0.9,0.9}{
            \makecell[{{p{\linewidth}}}]{
                \texttt{\tiny{[P1$\langle$GM]}}
                \texttt{Verplaatst 'M' naar (4,2) met succes. Uw bijgewerkte raster ziet er als volgt uit:} \\
\texttt{```} \\
\texttt{ 1234567} \\
\texttt{╔═╤══╤══╗} \\
\texttt{║T│◌◌│◌◌║ 1} \\
\texttt{╟─┤◌M│◌◌║ 2} \\
\texttt{║◌│PD├──╢ 3} \\
\texttt{╟─┘◌S│O◌║ 4} \\
\texttt{║◌◌◌◌│◌A║ 5} \\
\texttt{║◌◌┌─┼──╢ 6} \\
\texttt{║◌◌│◌│◌◌║ 7} \\
\texttt{╚══╧═╧══╝} \\
\texttt{```} \\
\texttt{Je speelt momenteel ronde 9 van maximaal 28.} \\
\texttt{Jullie hebben gezamenlijk 12 van 16 straffen verzameld.} \\
\texttt{De andere speler stuurde dit bericht:} \\
\\ 
\\ 
\texttt{"Ik heb S nu naar (3,4) verplaatst. Ik ga nu proberen de A naar de doelpositie (4,4) te verplaatsen."} \\
\\ 
\\ 
\texttt{Wat is je volgende opdracht?} \\
            }
        }
    }
    & \\ \\

    \theutterance \stepcounter{utterance}  
    & \multicolumn{4}{p{0.6\linewidth}}{
        \cellcolor[rgb]{0.9,0.9,0.9}{
            \makecell[{{p{\linewidth}}}]{
                \texttt{\tiny{[P1$\rangle$GM]}}
                \texttt{BEWEEG: M, (4,3)} \\
            }
        }
    }
    & & \\ \\

    \theutterance \stepcounter{utterance}  
    & & & \multicolumn{2}{p{0.3\linewidth}}{
        \cellcolor[rgb]{0.9,0.9,0.9}{
            \makecell[{{p{\linewidth}}}]{
                \texttt{\tiny{[GM$|$GM]}}
                \texttt{BEWEEG: M, (4,3)} \\
            }
        }
    }
    & & \\ \\

    \theutterance \stepcounter{utterance}  
    & & & \multicolumn{2}{p{0.3\linewidth}}{
        \cellcolor[rgb]{0.9,0.9,0.9}{
            \makecell[{{p{\linewidth}}}]{
                \texttt{\tiny{[GM$|$GM]}}
                \texttt{Straf: (4,3) is niet leeg, maar bevat 'D'.} \\
\\ 
\texttt{Jullie hebben gezamenlijk 13 van 16 straffen verzameld.} \\
\texttt{Probeer het opnieuw!} \\
            }
        }
    }
    & & \\ \\

    \theutterance \stepcounter{utterance}  
    & & \multicolumn{4}{p{0.6\linewidth}}{
        \cellcolor[rgb]{0.9,0.9,0.9}{
            \makecell[{{p{\linewidth}}}]{
                \texttt{\tiny{[P1$\langle$GM]}}
                \texttt{Straf: (4,3) is niet leeg, maar bevat 'D'.} \\
\\ 
\texttt{Jullie hebben gezamenlijk 13 van 16 straffen verzameld.} \\
\texttt{Probeer het opnieuw!} \\
            }
        }
    }
    & \\ \\

    \theutterance \stepcounter{utterance}  
    & \multicolumn{4}{p{0.6\linewidth}}{
        \cellcolor[rgb]{0.9,0.9,0.9}{
            \makecell[{{p{\linewidth}}}]{
                \texttt{\tiny{[P1$\rangle$GM]}}
                \texttt{BEWEEG: M, (4,1)} \\
            }
        }
    }
    & & \\ \\

    \theutterance \stepcounter{utterance}  
    & & & \multicolumn{2}{p{0.3\linewidth}}{
        \cellcolor[rgb]{0.9,0.9,0.9}{
            \makecell[{{p{\linewidth}}}]{
                \texttt{\tiny{[GM$|$GM]}}
                \texttt{BEWEEG: M, (4,1)} \\
            }
        }
    }
    & & \\ \\

    \theutterance \stepcounter{utterance}  
    & & & \multicolumn{2}{p{0.3\linewidth}}{
        \cellcolor[rgb]{0.9,0.9,0.9}{
            \makecell[{{p{\linewidth}}}]{
                \texttt{\tiny{[GM$|$GM]}}
                \texttt{M moved to (4, 1)} \\
            }
        }
    }
    & & \\ \\

    \theutterance \stepcounter{utterance}  
    & & \multicolumn{4}{p{0.6\linewidth}}{
        \cellcolor[rgb]{0.9,0.9,0.9}{
            \makecell[{{p{\linewidth}}}]{
                \texttt{\tiny{[GM$\rangle$P2]}}
                \texttt{Je bericht is doorgestuurd naar de andere speler.} \\
\\ 
\\ 
\texttt{Je speelt momenteel ronde 10 van maximaal 28.} \\
\texttt{Jullie hebben gezamenlijk 13 van 16 straffen verzameld.} \\
\texttt{De andere speler heeft een voorwerp op zijn rooster verplaatst.} \\
\\ 
\\ 
\texttt{Wat is je volgende opdracht?} \\
            }
        }
    }
    & \\ \\

    \theutterance \stepcounter{utterance}  
    & & & \multicolumn{4}{p{0.6\linewidth}}{
        \cellcolor[rgb]{0.9,0.9,0.9}{
            \makecell[{{p{\linewidth}}}]{
                \texttt{\tiny{[GM$\langle$P2]}}
                \texttt{BEWEEG: A, (4,4)} \\
            }
        }
    }
     \\ \\

    \theutterance \stepcounter{utterance}  
    & & & \multicolumn{2}{p{0.3\linewidth}}{
        \cellcolor[rgb]{0.9,0.9,0.9}{
            \makecell[{{p{\linewidth}}}]{
                \texttt{\tiny{[GM$|$GM]}}
                \texttt{BEWEEG: A, (4,4)} \\
            }
        }
    }
    & & \\ \\

    \theutterance \stepcounter{utterance}  
    & & & \multicolumn{2}{p{0.3\linewidth}}{
        \cellcolor[rgb]{0.9,0.9,0.9}{
            \makecell[{{p{\linewidth}}}]{
                \texttt{\tiny{[GM$|$GM]}}
                \texttt{A moved to (4, 4)} \\
            }
        }
    }
    & & \\ \\

    \theutterance \stepcounter{utterance}  
    & & \multicolumn{4}{p{0.6\linewidth}}{
        \cellcolor[rgb]{0.9,0.9,0.9}{
            \makecell[{{p{\linewidth}}}]{
                \texttt{\tiny{[P1$\langle$GM]}}
                \texttt{Verplaatst 'M' naar (4,1) met succes. Uw bijgewerkte raster ziet er als volgt uit:} \\
\texttt{```} \\
\texttt{ 1234567} \\
\texttt{╔═╤══╤══╗} \\
\texttt{║T│◌M│◌◌║ 1} \\
\texttt{╟─┤◌◌│◌◌║ 2} \\
\texttt{║◌│PD├──╢ 3} \\
\texttt{╟─┘◌S│O◌║ 4} \\
\texttt{║◌◌◌◌│◌A║ 5} \\
\texttt{║◌◌┌─┼──╢ 6} \\
\texttt{║◌◌│◌│◌◌║ 7} \\
\texttt{╚══╧═╧══╝} \\
\texttt{```} \\
\texttt{Je speelt momenteel ronde 10 van maximaal 28.} \\
\texttt{Jullie hebben gezamenlijk 13 van 16 straffen verzameld.} \\
\texttt{De andere speler heeft een voorwerp op zijn rooster verplaatst.} \\
\\ 
\\ 
\texttt{Wat is je volgende opdracht?} \\
            }
        }
    }
    & \\ \\

    \theutterance \stepcounter{utterance}  
    & \multicolumn{4}{p{0.6\linewidth}}{
        \cellcolor[rgb]{0.9,0.9,0.9}{
            \makecell[{{p{\linewidth}}}]{
                \texttt{\tiny{[P1$\rangle$GM]}}
                \texttt{ZEG: We hebben al 13 straffen, we moeten heel voorzichtig zijn! Ik heb M nu naar (4,1) verplaatst. Ik ga nu proberen om D naar de doelpositie (3,5) te verplaatsen.} \\
            }
        }
    }
    & & \\ \\

    \theutterance \stepcounter{utterance}  
    & & & \multicolumn{2}{p{0.3\linewidth}}{
        \cellcolor[rgb]{0.9,0.9,0.9}{
            \makecell[{{p{\linewidth}}}]{
                \texttt{\tiny{[GM$|$GM]}}
                \texttt{ZEG: We hebben al 13 straffen, we moeten heel voorzichtig zijn! Ik heb M nu naar (4,1) verplaatst. Ik ga nu proberen om D naar de doelpositie (3,5) te verplaatsen.} \\
            }
        }
    }
    & & \\ \\

    \theutterance \stepcounter{utterance}  
    & & \multicolumn{4}{p{0.6\linewidth}}{
        \cellcolor[rgb]{0.9,0.9,0.9}{
            \makecell[{{p{\linewidth}}}]{
                \texttt{\tiny{[GM$\rangle$P2]}}
                \texttt{Verplaatst 'A' naar (4,4) met succes. Uw bijgewerkte raster ziet er als volgt uit:} \\
\texttt{```} \\
\texttt{ 1234567} \\
\texttt{╔═╤══╤══╗} \\
\texttt{║◌│◌◌│◌T║ 1} \\
\texttt{╟─┤◌◌│◌◌║ 2} \\
\texttt{║◌│◌D├──╢ 3} \\
\texttt{╟─┘SA│◌◌║ 4} \\
\texttt{║◌◌◌O│◌◌║ 5} \\
\texttt{║PM┌─┼──╢ 6} \\
\texttt{║◌◌│◌│◌◌║ 7} \\
\texttt{╚══╧═╧══╝} \\
\texttt{```} \\
\texttt{Je speelt momenteel ronde 11 van maximaal 28.} \\
\texttt{Jullie hebben gezamenlijk 13 van 16 straffen verzameld.} \\
\texttt{De andere speler stuurde dit bericht:} \\
\\ 
\\ 
\texttt{"We hebben al 13 straffen, we moeten heel voorzichtig zijn! Ik heb M nu naar (4,1) verplaatst. Ik ga nu proberen om D naar de doelpositie (3,5) te verplaatsen."} \\
\\ 
\\ 
\texttt{Wat is je volgende opdracht?} \\
            }
        }
    }
    & \\ \\

    \theutterance \stepcounter{utterance}  
    & & & \multicolumn{4}{p{0.6\linewidth}}{
        \cellcolor[rgb]{0.9,0.9,0.9}{
            \makecell[{{p{\linewidth}}}]{
                \texttt{\tiny{[GM$\langle$P2]}}
                \texttt{ZEG: Je hebt gelijk, we moeten heel voorzichtig zijn met de straffen. Ik heb mijn A nu op (4,4) staan. Ik ga nu proberen mijn D van (4,3) naar de doelpositie (3,5) te verplaatsen.} \\
            }
        }
    }
     \\ \\

    \theutterance \stepcounter{utterance}  
    & & & \multicolumn{2}{p{0.3\linewidth}}{
        \cellcolor[rgb]{0.9,0.9,0.9}{
            \makecell[{{p{\linewidth}}}]{
                \texttt{\tiny{[GM$|$GM]}}
                \texttt{ZEG: Je hebt gelijk, we moeten heel voorzichtig zijn met de straffen. Ik heb mijn A nu op (4,4) staan. Ik ga nu proberen mijn D van (4,3) naar de doelpositie (3,5) te verplaatsen.} \\
            }
        }
    }
    & & \\ \\

    \theutterance \stepcounter{utterance}  
    & & \multicolumn{4}{p{0.6\linewidth}}{
        \cellcolor[rgb]{0.9,0.9,0.9}{
            \makecell[{{p{\linewidth}}}]{
                \texttt{\tiny{[P1$\langle$GM]}}
                \texttt{Je bericht is doorgestuurd naar de andere speler.} \\
\\ 
\\ 
\texttt{Je speelt momenteel ronde 11 van maximaal 28.} \\
\texttt{Jullie hebben gezamenlijk 13 van 16 straffen verzameld.} \\
\texttt{De andere speler stuurde dit bericht:} \\
\\ 
\\ 
\texttt{"Je hebt gelijk, we moeten heel voorzichtig zijn met de straffen. Ik heb mijn A nu op (4,4) staan. Ik ga nu proberen mijn D van (4,3) naar de doelpositie (3,5) te verplaatsen."} \\
\\ 
\\ 
\texttt{Wat is je volgende opdracht?} \\
            }
        }
    }
    & \\ \\

    \theutterance \stepcounter{utterance}  
    & \multicolumn{4}{p{0.6\linewidth}}{
        \cellcolor[rgb]{0.9,0.9,0.9}{
            \makecell[{{p{\linewidth}}}]{
                \texttt{\tiny{[P1$\rangle$GM]}}
                \texttt{BEWEEG: D, (3,5)} \\
            }
        }
    }
    & & \\ \\

    \theutterance \stepcounter{utterance}  
    & & & \multicolumn{2}{p{0.3\linewidth}}{
        \cellcolor[rgb]{0.9,0.9,0.9}{
            \makecell[{{p{\linewidth}}}]{
                \texttt{\tiny{[GM$|$GM]}}
                \texttt{BEWEEG: D, (3,5)} \\
            }
        }
    }
    & & \\ \\

    \theutterance \stepcounter{utterance}  
    & & & \multicolumn{2}{p{0.3\linewidth}}{
        \cellcolor[rgb]{0.9,0.9,0.9}{
            \makecell[{{p{\linewidth}}}]{
                \texttt{\tiny{[GM$|$GM]}}
                \texttt{D moved to (3, 5)} \\
            }
        }
    }
    & & \\ \\

    \theutterance \stepcounter{utterance}  
    & & \multicolumn{4}{p{0.6\linewidth}}{
        \cellcolor[rgb]{0.9,0.9,0.9}{
            \makecell[{{p{\linewidth}}}]{
                \texttt{\tiny{[GM$\rangle$P2]}}
                \texttt{Je bericht is doorgestuurd naar de andere speler.} \\
\\ 
\\ 
\texttt{Je speelt momenteel ronde 12 van maximaal 28.} \\
\texttt{Jullie hebben gezamenlijk 13 van 16 straffen verzameld.} \\
\texttt{De andere speler heeft een voorwerp op zijn rooster verplaatst.} \\
\\ 
\\ 
\texttt{Wat is je volgende opdracht?} \\
            }
        }
    }
    & \\ \\

    \theutterance \stepcounter{utterance}  
    & & & \multicolumn{4}{p{0.6\linewidth}}{
        \cellcolor[rgb]{0.9,0.9,0.9}{
            \makecell[{{p{\linewidth}}}]{
                \texttt{\tiny{[GM$\langle$P2]}}
                \texttt{BEWEEG: D, (3,5)} \\
            }
        }
    }
     \\ \\

    \theutterance \stepcounter{utterance}  
    & & & \multicolumn{2}{p{0.3\linewidth}}{
        \cellcolor[rgb]{0.9,0.9,0.9}{
            \makecell[{{p{\linewidth}}}]{
                \texttt{\tiny{[GM$|$GM]}}
                \texttt{BEWEEG: D, (3,5)} \\
            }
        }
    }
    & & \\ \\

    \theutterance \stepcounter{utterance}  
    & & & \multicolumn{2}{p{0.3\linewidth}}{
        \cellcolor[rgb]{0.9,0.9,0.9}{
            \makecell[{{p{\linewidth}}}]{
                \texttt{\tiny{[GM$|$GM]}}
                \texttt{D moved to (3, 5)} \\
            }
        }
    }
    & & \\ \\

    \theutterance \stepcounter{utterance}  
    & & \multicolumn{4}{p{0.6\linewidth}}{
        \cellcolor[rgb]{0.9,0.9,0.9}{
            \makecell[{{p{\linewidth}}}]{
                \texttt{\tiny{[P1$\langle$GM]}}
                \texttt{Verplaatst 'D' naar (3,5) met succes. Uw bijgewerkte raster ziet er als volgt uit:} \\
\texttt{```} \\
\texttt{ 1234567} \\
\texttt{╔═╤══╤══╗} \\
\texttt{║T│◌M│◌◌║ 1} \\
\texttt{╟─┤◌◌│◌◌║ 2} \\
\texttt{║◌│P◌├──╢ 3} \\
\texttt{╟─┘◌S│O◌║ 4} \\
\texttt{║◌◌D◌│◌A║ 5} \\
\texttt{║◌◌┌─┼──╢ 6} \\
\texttt{║◌◌│◌│◌◌║ 7} \\
\texttt{╚══╧═╧══╝} \\
\texttt{```} \\
\texttt{Je speelt momenteel ronde 12 van maximaal 28.} \\
\texttt{Jullie hebben gezamenlijk 13 van 16 straffen verzameld.} \\
\texttt{De andere speler heeft een voorwerp op zijn rooster verplaatst.} \\
\\ 
\\ 
\texttt{Wat is je volgende opdracht?} \\
            }
        }
    }
    & \\ \\

    \theutterance \stepcounter{utterance}  
    & \multicolumn{4}{p{0.6\linewidth}}{
        \cellcolor[rgb]{0.9,0.9,0.9}{
            \makecell[{{p{\linewidth}}}]{
                \texttt{\tiny{[P1$\rangle$GM]}}
                \texttt{ZEG: Goed gedaan! D staat nu op (3,5). Ik ga nu proberen om M naar de doelpositie (4,3) te verplaatsen.} \\
            }
        }
    }
    & & \\ \\

    \theutterance \stepcounter{utterance}  
    & & & \multicolumn{2}{p{0.3\linewidth}}{
        \cellcolor[rgb]{0.9,0.9,0.9}{
            \makecell[{{p{\linewidth}}}]{
                \texttt{\tiny{[GM$|$GM]}}
                \texttt{ZEG: Goed gedaan! D staat nu op (3,5). Ik ga nu proberen om M naar de doelpositie (4,3) te verplaatsen.} \\
            }
        }
    }
    & & \\ \\

    \theutterance \stepcounter{utterance}  
    & & \multicolumn{4}{p{0.6\linewidth}}{
        \cellcolor[rgb]{0.9,0.9,0.9}{
            \makecell[{{p{\linewidth}}}]{
                \texttt{\tiny{[GM$\rangle$P2]}}
                \texttt{Verplaatst 'D' naar (3,5) met succes. Uw bijgewerkte raster ziet er als volgt uit:} \\
\texttt{```} \\
\texttt{ 1234567} \\
\texttt{╔═╤══╤══╗} \\
\texttt{║◌│◌◌│◌T║ 1} \\
\texttt{╟─┤◌◌│◌◌║ 2} \\
\texttt{║◌│◌◌├──╢ 3} \\
\texttt{╟─┘SA│◌◌║ 4} \\
\texttt{║◌◌DO│◌◌║ 5} \\
\texttt{║PM┌─┼──╢ 6} \\
\texttt{║◌◌│◌│◌◌║ 7} \\
\texttt{╚══╧═╧══╝} \\
\texttt{```} \\
\texttt{Je speelt momenteel ronde 13 van maximaal 28.} \\
\texttt{Jullie hebben gezamenlijk 13 van 16 straffen verzameld.} \\
\texttt{De andere speler stuurde dit bericht:} \\
\\ 
\\ 
\texttt{"Goed gedaan! D staat nu op (3,5). Ik ga nu proberen om M naar de doelpositie (4,3) te verplaatsen."} \\
\\ 
\\ 
\texttt{Wat is je volgende opdracht?} \\
            }
        }
    }
    & \\ \\

    \theutterance \stepcounter{utterance}  
    & & & \multicolumn{4}{p{0.6\linewidth}}{
        \cellcolor[rgb]{0.9,0.9,0.9}{
            \makecell[{{p{\linewidth}}}]{
                \texttt{\tiny{[GM$\langle$P2]}}
                \texttt{ZEG: Goed bezig! Mijn D staat nu op (3,5). Ik ga nu proberen de P en M te verplaatsen naar hun doelposities (3,3) en (4,3). Ik begin met P naar (3,3).} \\
            }
        }
    }
     \\ \\

    \theutterance \stepcounter{utterance}  
    & & & \multicolumn{2}{p{0.3\linewidth}}{
        \cellcolor[rgb]{0.9,0.9,0.9}{
            \makecell[{{p{\linewidth}}}]{
                \texttt{\tiny{[GM$|$GM]}}
                \texttt{ZEG: Goed bezig! Mijn D staat nu op (3,5). Ik ga nu proberen de P en M te verplaatsen naar hun doelposities (3,3) en (4,3). Ik begin met P naar (3,3).} \\
            }
        }
    }
    & & \\ \\

    \theutterance \stepcounter{utterance}  
    & & \multicolumn{4}{p{0.6\linewidth}}{
        \cellcolor[rgb]{0.9,0.9,0.9}{
            \makecell[{{p{\linewidth}}}]{
                \texttt{\tiny{[P1$\langle$GM]}}
                \texttt{Je bericht is doorgestuurd naar de andere speler.} \\
\\ 
\\ 
\texttt{Je speelt momenteel ronde 13 van maximaal 28.} \\
\texttt{Jullie hebben gezamenlijk 13 van 16 straffen verzameld.} \\
\texttt{De andere speler stuurde dit bericht:} \\
\\ 
\\ 
\texttt{"Goed bezig! Mijn D staat nu op (3,5). Ik ga nu proberen de P en M te verplaatsen naar hun doelposities (3,3) en (4,3). Ik begin met P naar (3,3)."} \\
\\ 
\\ 
\texttt{Wat is je volgende opdracht?} \\
            }
        }
    }
    & \\ \\

    \theutterance \stepcounter{utterance}  
    & \multicolumn{4}{p{0.6\linewidth}}{
        \cellcolor[rgb]{0.9,0.9,0.9}{
            \makecell[{{p{\linewidth}}}]{
                \texttt{\tiny{[P1$\rangle$GM]}}
                \texttt{BEWEEG: M, (4,3)} \\
            }
        }
    }
    & & \\ \\

    \theutterance \stepcounter{utterance}  
    & & & \multicolumn{2}{p{0.3\linewidth}}{
        \cellcolor[rgb]{0.9,0.9,0.9}{
            \makecell[{{p{\linewidth}}}]{
                \texttt{\tiny{[GM$|$GM]}}
                \texttt{BEWEEG: M, (4,3)} \\
            }
        }
    }
    & & \\ \\

    \theutterance \stepcounter{utterance}  
    & & & \multicolumn{2}{p{0.3\linewidth}}{
        \cellcolor[rgb]{0.9,0.9,0.9}{
            \makecell[{{p{\linewidth}}}]{
                \texttt{\tiny{[GM$|$GM]}}
                \texttt{M moved to (4, 3)} \\
            }
        }
    }
    & & \\ \\

    \theutterance \stepcounter{utterance}  
    & & \multicolumn{4}{p{0.6\linewidth}}{
        \cellcolor[rgb]{0.9,0.9,0.9}{
            \makecell[{{p{\linewidth}}}]{
                \texttt{\tiny{[GM$\rangle$P2]}}
                \texttt{Je bericht is doorgestuurd naar de andere speler.} \\
\\ 
\\ 
\texttt{Je speelt momenteel ronde 14 van maximaal 28.} \\
\texttt{Jullie hebben gezamenlijk 13 van 16 straffen verzameld.} \\
\texttt{De andere speler heeft een voorwerp op zijn rooster verplaatst.} \\
\\ 
\\ 
\texttt{Wat is je volgende opdracht?} \\
            }
        }
    }
    & \\ \\

    \theutterance \stepcounter{utterance}  
    & & & \multicolumn{4}{p{0.6\linewidth}}{
        \cellcolor[rgb]{0.9,0.9,0.9}{
            \makecell[{{p{\linewidth}}}]{
                \texttt{\tiny{[GM$\langle$P2]}}
                \texttt{BEWEEG: P, (3,3)} \\
            }
        }
    }
     \\ \\

    \theutterance \stepcounter{utterance}  
    & & & \multicolumn{2}{p{0.3\linewidth}}{
        \cellcolor[rgb]{0.9,0.9,0.9}{
            \makecell[{{p{\linewidth}}}]{
                \texttt{\tiny{[GM$|$GM]}}
                \texttt{BEWEEG: P, (3,3)} \\
            }
        }
    }
    & & \\ \\

    \theutterance \stepcounter{utterance}  
    & & & \multicolumn{2}{p{0.3\linewidth}}{
        \cellcolor[rgb]{0.9,0.9,0.9}{
            \makecell[{{p{\linewidth}}}]{
                \texttt{\tiny{[GM$|$GM]}}
                \texttt{P moved to (3, 3)} \\
            }
        }
    }
    & & \\ \\

    \theutterance \stepcounter{utterance}  
    & & \multicolumn{4}{p{0.6\linewidth}}{
        \cellcolor[rgb]{0.9,0.9,0.9}{
            \makecell[{{p{\linewidth}}}]{
                \texttt{\tiny{[P1$\langle$GM]}}
                \texttt{Verplaatst 'M' naar (4,3) met succes. Uw bijgewerkte raster ziet er als volgt uit:} \\
\texttt{```} \\
\texttt{ 1234567} \\
\texttt{╔═╤══╤══╗} \\
\texttt{║T│◌◌│◌◌║ 1} \\
\texttt{╟─┤◌◌│◌◌║ 2} \\
\texttt{║◌│PM├──╢ 3} \\
\texttt{╟─┘◌S│O◌║ 4} \\
\texttt{║◌◌D◌│◌A║ 5} \\
\texttt{║◌◌┌─┼──╢ 6} \\
\texttt{║◌◌│◌│◌◌║ 7} \\
\texttt{╚══╧═╧══╝} \\
\texttt{```} \\
\texttt{Je speelt momenteel ronde 14 van maximaal 28.} \\
\texttt{Jullie hebben gezamenlijk 13 van 16 straffen verzameld.} \\
\texttt{De andere speler heeft een voorwerp op zijn rooster verplaatst.} \\
\\ 
\\ 
\texttt{Wat is je volgende opdracht?} \\
            }
        }
    }
    & \\ \\

    \theutterance \stepcounter{utterance}  
    & \multicolumn{4}{p{0.6\linewidth}}{
        \cellcolor[rgb]{0.9,0.9,0.9}{
            \makecell[{{p{\linewidth}}}]{
                \texttt{\tiny{[P1$\rangle$GM]}}
                \texttt{ZEG: Geweldig, M staat nu op (4,3)! Dat is een doelpositie. Ik heb P op (3,3) en M op (4,3), D op (3,5), S op (4,4), O op (6,4), A op (7,5) en T op (1,1). We moeten nog T naar (5,5) en O naar (4,5) en A naar (4,4) proberen te krijgen. Wacht, ik zie dat ik A en S nog moet checken. Laten we kijken wat de status is.} \\
            }
        }
    }
    & & \\ \\

    \theutterance \stepcounter{utterance}  
    & & & \multicolumn{2}{p{0.3\linewidth}}{
        \cellcolor[rgb]{0.9,0.9,0.9}{
            \makecell[{{p{\linewidth}}}]{
                \texttt{\tiny{[GM$|$GM]}}
                \texttt{ZEG: Geweldig, M staat nu op (4,3)! Dat is een doelpositie. Ik heb P op (3,3) en M op (4,3), D op (3,5), S op (4,4), O op (6,4), A op (7,5) en T op (1,1). We moeten nog T naar (5,5) en O naar (4,5) en A naar (4,4) proberen te krijgen. Wacht, ik zie dat ik A en S nog moet checken. Laten we kijken wat de status is.} \\
            }
        }
    }
    & & \\ \\

    \theutterance \stepcounter{utterance}  
    & & \multicolumn{4}{p{0.6\linewidth}}{
        \cellcolor[rgb]{0.9,0.9,0.9}{
            \makecell[{{p{\linewidth}}}]{
                \texttt{\tiny{[GM$\rangle$P2]}}
                \texttt{Verplaatst 'P' naar (3,3) met succes. Uw bijgewerkte raster ziet er als volgt uit:} \\
\texttt{```} \\
\texttt{ 1234567} \\
\texttt{╔═╤══╤══╗} \\
\texttt{║◌│◌◌│◌T║ 1} \\
\texttt{╟─┤◌◌│◌◌║ 2} \\
\texttt{║◌│P◌├──╢ 3} \\
\texttt{╟─┘SA│◌◌║ 4} \\
\texttt{║◌◌DO│◌◌║ 5} \\
\texttt{║◌M┌─┼──╢ 6} \\
\texttt{║◌◌│◌│◌◌║ 7} \\
\texttt{╚══╧═╧══╝} \\
\texttt{```} \\
\texttt{Je speelt momenteel ronde 15 van maximaal 28.} \\
\texttt{Jullie hebben gezamenlijk 13 van 16 straffen verzameld.} \\
\texttt{De andere speler stuurde dit bericht:} \\
\\ 
\\ 
\texttt{"Geweldig, M staat nu op (4,3)! Dat is een doelpositie. Ik heb P op (3,3) en M op (4,3), D op (3,5), S op (4,4), O op (6,4), A op (7,5) en T op (1,1). We moeten nog T naar (5,5) en O naar (4,5) en A naar (4,4) proberen te krijgen. Wacht, ik zie dat ik A en S nog moet checken. Laten we kijken wat de status is."} \\
\\ 
\\ 
\texttt{Wat is je volgende opdracht?} \\
            }
        }
    }
    & \\ \\

    \theutterance \stepcounter{utterance}  
    & & & \multicolumn{4}{p{0.6\linewidth}}{
        \cellcolor[rgb]{0.9,0.9,0.9}{
            \makecell[{{p{\linewidth}}}]{
                \texttt{\tiny{[GM$\langle$P2]}}
                \texttt{ZEG: Ik heb P op (3,3), M op (1,6) {-} wacht, mijn M staat nog op (2,6) volgens het raster. Mijn huidige posities zijn: P(3,3), M(2,6), S(3,4), A(4,4), D(3,5), O(4,5) en T(7,1). Laten we de laatste objecten doen: ik ga nu M naar (4,3) verplaatsen.} \\
            }
        }
    }
     \\ \\

    \theutterance \stepcounter{utterance}  
    & & & \multicolumn{2}{p{0.3\linewidth}}{
        \cellcolor[rgb]{0.9,0.9,0.9}{
            \makecell[{{p{\linewidth}}}]{
                \texttt{\tiny{[GM$|$GM]}}
                \texttt{ZEG: Ik heb P op (3,3), M op (1,6) {-} wacht, mijn M staat nog op (2,6) volgens het raster. Mijn huidige posities zijn: P(3,3), M(2,6), S(3,4), A(4,4), D(3,5), O(4,5) en T(7,1). Laten we de laatste objecten doen: ik ga nu M naar (4,3) verplaatsen.} \\
            }
        }
    }
    & & \\ \\

    \theutterance \stepcounter{utterance}  
    & & \multicolumn{4}{p{0.6\linewidth}}{
        \cellcolor[rgb]{0.9,0.9,0.9}{
            \makecell[{{p{\linewidth}}}]{
                \texttt{\tiny{[P1$\langle$GM]}}
                \texttt{Je bericht is doorgestuurd naar de andere speler.} \\
\\ 
\\ 
\texttt{Je speelt momenteel ronde 15 van maximaal 28.} \\
\texttt{Jullie hebben gezamenlijk 13 van 16 straffen verzameld.} \\
\texttt{De andere speler stuurde dit bericht:} \\
\\ 
\\ 
\texttt{"Ik heb P op (3,3), M op (1,6) {-} wacht, mijn M staat nog op (2,6) volgens het raster. Mijn huidige posities zijn: P(3,3), M(2,6), S(3,4), A(4,4), D(3,5), O(4,5) en T(7,1). Laten we de laatste objecten doen: ik ga nu M naar (4,3) verplaatsen."} \\
\\ 
\\ 
\texttt{Wat is je volgende opdracht?} \\
            }
        }
    }
    & \\ \\

    \theutterance \stepcounter{utterance}  
    & \multicolumn{4}{p{0.6\linewidth}}{
        \cellcolor[rgb]{0.9,0.9,0.9}{
            \makecell[{{p{\linewidth}}}]{
                \texttt{\tiny{[P1$\rangle$GM]}}
                \texttt{ZEG: Dat is een goed plan. Ik heb mijn objecten bijna op de juiste plek. Mijn huidige posities zijn: T(1,1), P(3,3), M(4,3), S(4,4), D(3,5), O(6,4) en A(7,5). Ik ga nu proberen om mijn O naar de doelpositie (4,5) te verplaatsen.} \\
            }
        }
    }
    & & \\ \\

    \theutterance \stepcounter{utterance}  
    & & & \multicolumn{2}{p{0.3\linewidth}}{
        \cellcolor[rgb]{0.9,0.9,0.9}{
            \makecell[{{p{\linewidth}}}]{
                \texttt{\tiny{[GM$|$GM]}}
                \texttt{ZEG: Dat is een goed plan. Ik heb mijn objecten bijna op de juiste plek. Mijn huidige posities zijn: T(1,1), P(3,3), M(4,3), S(4,4), D(3,5), O(6,4) en A(7,5). Ik ga nu proberen om mijn O naar de doelpositie (4,5) te verplaatsen.} \\
            }
        }
    }
    & & \\ \\

    \theutterance \stepcounter{utterance}  
    & & \multicolumn{4}{p{0.6\linewidth}}{
        \cellcolor[rgb]{0.9,0.9,0.9}{
            \makecell[{{p{\linewidth}}}]{
                \texttt{\tiny{[GM$\rangle$P2]}}
                \texttt{Je bericht is doorgestuurd naar de andere speler.} \\
\\ 
\\ 
\texttt{Je speelt momenteel ronde 16 van maximaal 28.} \\
\texttt{Jullie hebben gezamenlijk 13 van 16 straffen verzameld.} \\
\texttt{De andere speler stuurde dit bericht:} \\
\\ 
\\ 
\texttt{"Dat is een goed plan. Ik heb mijn objecten bijna op de juiste plek. Mijn huidige posities zijn: T(1,1), P(3,3), M(4,3), S(4,4), D(3,5), O(6,4) en A(7,5). Ik ga nu proberen om mijn O naar de doelpositie (4,5) te verplaatsen."} \\
\\ 
\\ 
\texttt{Wat is je volgende opdracht?} \\
            }
        }
    }
    & \\ \\

    \theutterance \stepcounter{utterance}  
    & & & \multicolumn{4}{p{0.6\linewidth}}{
        \cellcolor[rgb]{0.9,0.9,0.9}{
            \makecell[{{p{\linewidth}}}]{
                \texttt{\tiny{[GM$\langle$P2]}}
                \texttt{BEWEEG: M, (4,3)} \\
            }
        }
    }
     \\ \\

    \theutterance \stepcounter{utterance}  
    & & & \multicolumn{2}{p{0.3\linewidth}}{
        \cellcolor[rgb]{0.9,0.9,0.9}{
            \makecell[{{p{\linewidth}}}]{
                \texttt{\tiny{[GM$|$GM]}}
                \texttt{BEWEEG: M, (4,3)} \\
            }
        }
    }
    & & \\ \\

    \theutterance \stepcounter{utterance}  
    & & & \multicolumn{2}{p{0.3\linewidth}}{
        \cellcolor[rgb]{0.9,0.9,0.9}{
            \makecell[{{p{\linewidth}}}]{
                \texttt{\tiny{[GM$|$GM]}}
                \texttt{M moved to (4, 3)} \\
            }
        }
    }
    & & \\ \\

    \theutterance \stepcounter{utterance}  
    & & \multicolumn{4}{p{0.6\linewidth}}{
        \cellcolor[rgb]{0.9,0.9,0.9}{
            \makecell[{{p{\linewidth}}}]{
                \texttt{\tiny{[P1$\langle$GM]}}
                \texttt{Je bericht is doorgestuurd naar de andere speler.} \\
\\ 
\\ 
\texttt{Je speelt momenteel ronde 16 van maximaal 28.} \\
\texttt{Jullie hebben gezamenlijk 13 van 16 straffen verzameld.} \\
\texttt{De andere speler heeft een voorwerp op zijn rooster verplaatst.} \\
\\ 
\\ 
\texttt{Wat is je volgende opdracht?} \\
            }
        }
    }
    & \\ \\

    \theutterance \stepcounter{utterance}  
    & \multicolumn{4}{p{0.6\linewidth}}{
        \cellcolor[rgb]{0.9,0.9,0.9}{
            \makecell[{{p{\linewidth}}}]{
                \texttt{\tiny{[P1$\rangle$GM]}}
                \texttt{BEWEEG: O, (4,5)} \\
            }
        }
    }
    & & \\ \\

    \theutterance \stepcounter{utterance}  
    & & & \multicolumn{2}{p{0.3\linewidth}}{
        \cellcolor[rgb]{0.9,0.9,0.9}{
            \makecell[{{p{\linewidth}}}]{
                \texttt{\tiny{[GM$|$GM]}}
                \texttt{BEWEEG: O, (4,5)} \\
            }
        }
    }
    & & \\ \\

    \theutterance \stepcounter{utterance}  
    & & & \multicolumn{2}{p{0.3\linewidth}}{
        \cellcolor[rgb]{0.9,0.9,0.9}{
            \makecell[{{p{\linewidth}}}]{
                \texttt{\tiny{[GM$|$GM]}}
                \texttt{O moved to (4, 5)} \\
            }
        }
    }
    & & \\ \\

    \theutterance \stepcounter{utterance}  
    & & \multicolumn{4}{p{0.6\linewidth}}{
        \cellcolor[rgb]{0.9,0.9,0.9}{
            \makecell[{{p{\linewidth}}}]{
                \texttt{\tiny{[GM$\rangle$P2]}}
                \texttt{Verplaatst 'M' naar (4,3) met succes. Uw bijgewerkte raster ziet er als volgt uit:} \\
\texttt{```} \\
\texttt{ 1234567} \\
\texttt{╔═╤══╤══╗} \\
\texttt{║◌│◌◌│◌T║ 1} \\
\texttt{╟─┤◌◌│◌◌║ 2} \\
\texttt{║◌│PM├──╢ 3} \\
\texttt{╟─┘SA│◌◌║ 4} \\
\texttt{║◌◌DO│◌◌║ 5} \\
\texttt{║◌◌┌─┼──╢ 6} \\
\texttt{║◌◌│◌│◌◌║ 7} \\
\texttt{╚══╧═╧══╝} \\
\texttt{```} \\
\texttt{Je speelt momenteel ronde 17 van maximaal 28.} \\
\texttt{Jullie hebben gezamenlijk 13 van 16 straffen verzameld.} \\
\texttt{De andere speler heeft een voorwerp op zijn rooster verplaatst.} \\
\\ 
\\ 
\texttt{Wat is je volgende opdracht?} \\
            }
        }
    }
    & \\ \\

    \theutterance \stepcounter{utterance}  
    & & & \multicolumn{4}{p{0.6\linewidth}}{
        \cellcolor[rgb]{0.9,0.9,0.9}{
            \makecell[{{p{\linewidth}}}]{
                \texttt{\tiny{[GM$\langle$P2]}}
                \texttt{ZEG: Ik heb mijn M nu op (4,3) staan. Mijn huidige posities zijn: P(3,3), M(4,3), S(3,4), A(4,4), D(3,5), O(4,5) en T(7,1). We moeten alleen nog T naar (5,5) verplaatsen om klaar te zijn!} \\
            }
        }
    }
     \\ \\

    \theutterance \stepcounter{utterance}  
    & & & \multicolumn{2}{p{0.3\linewidth}}{
        \cellcolor[rgb]{0.9,0.9,0.9}{
            \makecell[{{p{\linewidth}}}]{
                \texttt{\tiny{[GM$|$GM]}}
                \texttt{ZEG: Ik heb mijn M nu op (4,3) staan. Mijn huidige posities zijn: P(3,3), M(4,3), S(3,4), A(4,4), D(3,5), O(4,5) en T(7,1). We moeten alleen nog T naar (5,5) verplaatsen om klaar te zijn!} \\
            }
        }
    }
    & & \\ \\

    \theutterance \stepcounter{utterance}  
    & & \multicolumn{4}{p{0.6\linewidth}}{
        \cellcolor[rgb]{0.9,0.9,0.9}{
            \makecell[{{p{\linewidth}}}]{
                \texttt{\tiny{[P1$\langle$GM]}}
                \texttt{Verplaatst 'O' naar (4,5) met succes. Uw bijgewerkte raster ziet er als volgt uit:} \\
\texttt{```} \\
\texttt{ 1234567} \\
\texttt{╔═╤══╤══╗} \\
\texttt{║T│◌◌│◌◌║ 1} \\
\texttt{╟─┤◌◌│◌◌║ 2} \\
\texttt{║◌│PM├──╢ 3} \\
\texttt{╟─┘◌S│◌◌║ 4} \\
\texttt{║◌◌DO│◌A║ 5} \\
\texttt{║◌◌┌─┼──╢ 6} \\
\texttt{║◌◌│◌│◌◌║ 7} \\
\texttt{╚══╧═╧══╝} \\
\texttt{```} \\
\texttt{Je speelt momenteel ronde 17 van maximaal 28.} \\
\texttt{Jullie hebben gezamenlijk 13 van 16 straffen verzameld.} \\
\texttt{De andere speler stuurde dit bericht:} \\
\\ 
\\ 
\texttt{"Ik heb mijn M nu op (4,3) staan. Mijn huidige posities zijn: P(3,3), M(4,3), S(3,4), A(4,4), D(3,5), O(4,5) en T(7,1). We moeten alleen nog T naar (5,5) verplaatsen om klaar te zijn!"} \\
\\ 
\\ 
\texttt{Wat is je volgende opdracht?} \\
            }
        }
    }
    & \\ \\

    \theutterance \stepcounter{utterance}  
    & \multicolumn{4}{p{0.6\linewidth}}{
        \cellcolor[rgb]{0.9,0.9,0.9}{
            \makecell[{{p{\linewidth}}}]{
                \texttt{\tiny{[P1$\rangle$GM]}}
                \texttt{ZEG: Dat is fantastisch! We zijn er bijna. Mijn posities zijn nu: T(1,1), P(3,3), M(4,3), S(4,4), D(3,5), O(4,5), A(7,5). Ik moet alleen nog mijn T naar (5,5) en mijn A naar (4,4) verplaatsen. Maar wacht, mijn S staat al op (4,4), dus ik moet eerst S verplaatsen of kijken of de andere speler de A al op (4,4) heeft. Laten we eerst mijn T naar (5,5) proberen te krijgen.} \\
            }
        }
    }
    & & \\ \\

    \theutterance \stepcounter{utterance}  
    & & & \multicolumn{2}{p{0.3\linewidth}}{
        \cellcolor[rgb]{0.9,0.9,0.9}{
            \makecell[{{p{\linewidth}}}]{
                \texttt{\tiny{[GM$|$GM]}}
                \texttt{ZEG: Dat is fantastisch! We zijn er bijna. Mijn posities zijn nu: T(1,1), P(3,3), M(4,3), S(4,4), D(3,5), O(4,5), A(7,5). Ik moet alleen nog mijn T naar (5,5) en mijn A naar (4,4) verplaatsen. Maar wacht, mijn S staat al op (4,4), dus ik moet eerst S verplaatsen of kijken of de andere speler de A al op (4,4) heeft. Laten we eerst mijn T naar (5,5) proberen te krijgen.} \\
            }
        }
    }
    & & \\ \\

    \theutterance \stepcounter{utterance}  
    & & \multicolumn{4}{p{0.6\linewidth}}{
        \cellcolor[rgb]{0.9,0.9,0.9}{
            \makecell[{{p{\linewidth}}}]{
                \texttt{\tiny{[GM$\rangle$P2]}}
                \texttt{Je bericht is doorgestuurd naar de andere speler.} \\
\\ 
\\ 
\texttt{Je speelt momenteel ronde 18 van maximaal 28.} \\
\texttt{Jullie hebben gezamenlijk 13 van 16 straffen verzameld.} \\
\texttt{De andere speler stuurde dit bericht:} \\
\\ 
\\ 
\texttt{"Dat is fantastisch! We zijn er bijna. Mijn posities zijn nu: T(1,1), P(3,3), M(4,3), S(4,4), D(3,5), O(4,5), A(7,5). Ik moet alleen nog mijn T naar (5,5) en mijn A naar (4,4) verplaatsen. Maar wacht, mijn S staat al op (4,4), dus ik moet eerst S verplaatsen of kijken of de andere speler de A al op (4,4) heeft. Laten we eerst mijn T naar (5,5) proberen te krijgen."} \\
\\ 
\\ 
\texttt{Wat is je volgende opdracht?} \\
            }
        }
    }
    & \\ \\

    \theutterance \stepcounter{utterance}  
    & & & \multicolumn{4}{p{0.6\linewidth}}{
        \cellcolor[rgb]{0.9,0.9,0.9}{
            \makecell[{{p{\linewidth}}}]{
                \texttt{\tiny{[GM$\langle$P2]}}
                \texttt{BEWEEG: T, (5,5)} \\
            }
        }
    }
     \\ \\

    \theutterance \stepcounter{utterance}  
    & & & \multicolumn{2}{p{0.3\linewidth}}{
        \cellcolor[rgb]{0.9,0.9,0.9}{
            \makecell[{{p{\linewidth}}}]{
                \texttt{\tiny{[GM$|$GM]}}
                \texttt{BEWEEG: T, (5,5)} \\
            }
        }
    }
    & & \\ \\

    \theutterance \stepcounter{utterance}  
    & & & \multicolumn{2}{p{0.3\linewidth}}{
        \cellcolor[rgb]{0.9,0.9,0.9}{
            \makecell[{{p{\linewidth}}}]{
                \texttt{\tiny{[GM$|$GM]}}
                \texttt{Straf: (5,5) is niet leeg, maar bevat '│'.} \\
\\ 
\texttt{Jullie hebben gezamenlijk 14 van 16 straffen verzameld.} \\
\texttt{Probeer het opnieuw!} \\
            }
        }
    }
    & & \\ \\

    \theutterance \stepcounter{utterance}  
    & & \multicolumn{4}{p{0.6\linewidth}}{
        \cellcolor[rgb]{0.9,0.9,0.9}{
            \makecell[{{p{\linewidth}}}]{
                \texttt{\tiny{[GM$\rangle$P2]}}
                \texttt{Straf: (5,5) is niet leeg, maar bevat '│'.} \\
\\ 
\texttt{Jullie hebben gezamenlijk 14 van 16 straffen verzameld.} \\
\texttt{Probeer het opnieuw!} \\
            }
        }
    }
    & \\ \\

    \theutterance \stepcounter{utterance}  
    & & & \multicolumn{4}{p{0.6\linewidth}}{
        \cellcolor[rgb]{0.9,0.9,0.9}{
            \makecell[{{p{\linewidth}}}]{
                \texttt{\tiny{[GM$\langle$P2]}}
                \texttt{BEWEEG: T, (6,5)} \\
            }
        }
    }
     \\ \\

    \theutterance \stepcounter{utterance}  
    & & & \multicolumn{2}{p{0.3\linewidth}}{
        \cellcolor[rgb]{0.9,0.9,0.9}{
            \makecell[{{p{\linewidth}}}]{
                \texttt{\tiny{[GM$|$GM]}}
                \texttt{BEWEEG: T, (6,5)} \\
            }
        }
    }
    & & \\ \\

    \theutterance \stepcounter{utterance}  
    & & & \multicolumn{2}{p{0.3\linewidth}}{
        \cellcolor[rgb]{0.9,0.9,0.9}{
            \makecell[{{p{\linewidth}}}]{
                \texttt{\tiny{[GM$|$GM]}}
                \texttt{T moved to (6, 5)} \\
            }
        }
    }
    & & \\ \\

    \theutterance \stepcounter{utterance}  
    & & \multicolumn{4}{p{0.6\linewidth}}{
        \cellcolor[rgb]{0.9,0.9,0.9}{
            \makecell[{{p{\linewidth}}}]{
                \texttt{\tiny{[P1$\langle$GM]}}
                \texttt{Je bericht is doorgestuurd naar de andere speler.} \\
\\ 
\\ 
\texttt{Je speelt momenteel ronde 18 van maximaal 28.} \\
\texttt{Jullie hebben gezamenlijk 14 van 16 straffen verzameld.} \\
\texttt{De andere speler heeft een voorwerp op zijn rooster verplaatst.} \\
\\ 
\\ 
\texttt{Wat is je volgende opdracht?} \\
            }
        }
    }
    & \\ \\

    \theutterance \stepcounter{utterance}  
    & \multicolumn{4}{p{0.6\linewidth}}{
        \cellcolor[rgb]{0.9,0.9,0.9}{
            \makecell[{{p{\linewidth}}}]{
                \texttt{\tiny{[P1$\rangle$GM]}}
                \texttt{BEWEEG: T, (5,5)} \\
            }
        }
    }
    & & \\ \\

    \theutterance \stepcounter{utterance}  
    & & & \multicolumn{2}{p{0.3\linewidth}}{
        \cellcolor[rgb]{0.9,0.9,0.9}{
            \makecell[{{p{\linewidth}}}]{
                \texttt{\tiny{[GM$|$GM]}}
                \texttt{BEWEEG: T, (5,5)} \\
            }
        }
    }
    & & \\ \\

    \theutterance \stepcounter{utterance}  
    & & & \multicolumn{2}{p{0.3\linewidth}}{
        \cellcolor[rgb]{0.9,0.9,0.9}{
            \makecell[{{p{\linewidth}}}]{
                \texttt{\tiny{[GM$|$GM]}}
                \texttt{Straf: (5,5) is niet leeg, maar bevat '│'.} \\
\\ 
\texttt{Jullie hebben gezamenlijk 15 van 16 straffen verzameld.} \\
\texttt{Probeer het opnieuw!} \\
            }
        }
    }
    & & \\ \\

    \theutterance \stepcounter{utterance}  
    & & \multicolumn{4}{p{0.6\linewidth}}{
        \cellcolor[rgb]{0.9,0.9,0.9}{
            \makecell[{{p{\linewidth}}}]{
                \texttt{\tiny{[P1$\langle$GM]}}
                \texttt{Straf: (5,5) is niet leeg, maar bevat '│'.} \\
\\ 
\texttt{Jullie hebben gezamenlijk 15 van 16 straffen verzameld.} \\
\texttt{Probeer het opnieuw!} \\
            }
        }
    }
    & \\ \\

    \theutterance \stepcounter{utterance}  
    & \multicolumn{4}{p{0.6\linewidth}}{
        \cellcolor[rgb]{0.9,0.9,0.9}{
            \makecell[{{p{\linewidth}}}]{
                \texttt{\tiny{[P1$\rangle$GM]}}
                \texttt{BEWEEG: T, (5,1)} \\
            }
        }
    }
    & & \\ \\

    \theutterance \stepcounter{utterance}  
    & & & \multicolumn{2}{p{0.3\linewidth}}{
        \cellcolor[rgb]{0.9,0.9,0.9}{
            \makecell[{{p{\linewidth}}}]{
                \texttt{\tiny{[GM$|$GM]}}
                \texttt{BEWEEG: T, (5,1)} \\
            }
        }
    }
    & & \\ \\

    \theutterance \stepcounter{utterance}  
    & & & \multicolumn{2}{p{0.3\linewidth}}{
        \cellcolor[rgb]{0.9,0.9,0.9}{
            \makecell[{{p{\linewidth}}}]{
                \texttt{\tiny{[GM$|$GM]}}
                \texttt{Straf: (5,1) is niet leeg, maar bevat '│'.} \\
\\ 
\texttt{Jullie hebben gezamenlijk 16 van 16 straffen verzameld.} \\
\texttt{Probeer het opnieuw!} \\
            }
        }
    }
    & & \\ \\

    \theutterance \stepcounter{utterance}  
    & & \multicolumn{4}{p{0.6\linewidth}}{
        \cellcolor[rgb]{0.9,0.9,0.9}{
            \makecell[{{p{\linewidth}}}]{
                \texttt{\tiny{[P1$\langle$GM]}}
                \texttt{Straf: (5,1) is niet leeg, maar bevat '│'.} \\
\\ 
\texttt{Jullie hebben gezamenlijk 16 van 16 straffen verzameld.} \\
\texttt{Probeer het opnieuw!} \\
            }
        }
    }
    & \\ \\

    \theutterance \stepcounter{utterance}  
    & \multicolumn{4}{p{0.6\linewidth}}{
        \cellcolor[rgb]{0.9,0.9,0.9}{
            \makecell[{{p{\linewidth}}}]{
                \texttt{\tiny{[P1$\rangle$GM]}}
                \texttt{BEWEEG: T, (5,2)} \\
            }
        }
    }
    & & \\ \\

    \theutterance \stepcounter{utterance}  
    & & & \multicolumn{2}{p{0.3\linewidth}}{
        \cellcolor[rgb]{0.9,0.9,0.9}{
            \makecell[{{p{\linewidth}}}]{
                \texttt{\tiny{[GM$|$GM]}}
                \texttt{BEWEEG: T, (5,2)} \\
            }
        }
    }
    & & \\ \\

    \theutterance \stepcounter{utterance}  
    & & & \multicolumn{2}{p{0.3\linewidth}}{
        \cellcolor[rgb]{0.9,0.9,0.9}{
            \makecell[{{p{\linewidth}}}]{
                \texttt{\tiny{[GM$|$GM]}}
                \texttt{Straf: (5,2) is niet leeg, maar bevat '│'.} \\
\\ 
\texttt{Jullie hebben gezamenlijk 17 van 16 straffen verzameld.} \\
\texttt{Probeer het opnieuw!} \\
            }
        }
    }
    & & \\ \\

    \theutterance \stepcounter{utterance}  
    & & \multicolumn{4}{p{0.6\linewidth}}{
        \cellcolor[rgb]{0.9,0.9,0.9}{
            \makecell[{{p{\linewidth}}}]{
                \texttt{\tiny{[P1$\langle$GM]}}
                \texttt{Straf: (5,2) is niet leeg, maar bevat '│'.} \\
\\ 
\texttt{Jullie hebben gezamenlijk 17 van 16 straffen verzameld.} \\
\texttt{Probeer het opnieuw!} \\
            }
        }
    }
    & \\ \\

    \theutterance \stepcounter{utterance}  
    & \multicolumn{4}{p{0.6\linewidth}}{
        \cellcolor[rgb]{0.9,0.9,0.9}{
            \makecell[{{p{\linewidth}}}]{
                \texttt{\tiny{[P1$\rangle$GM]}}
                \texttt{Ik heb de regels verkeerd geïnterpreteerd door te proberen te bewegen naar cellen die onderdeel zijn van het kader (de lijnen). Ik heb nu 17 strafpunten verzameld, wat betekent dat we het spel hebben verloren. Ik bied mijn excuses aan voor de foutieve bewegingen.} \\
            }
        }
    }
    & & \\ \\

    \theutterance \stepcounter{utterance}  
    & & & \multicolumn{2}{p{0.3\linewidth}}{
        \cellcolor[rgb]{0.9,0.9,0.9}{
            \makecell[{{p{\linewidth}}}]{
                \texttt{\tiny{[GM$|$GM]}}
                \texttt{Ik heb de regels verkeerd geïnterpreteerd door te proberen te bewegen naar cellen die onderdeel zijn van het kader (de lijnen). Ik heb nu 17 strafpunten verzameld, wat betekent dat we het spel hebben verloren. Ik bied mijn excuses aan voor de foutieve bewegingen.} \\
            }
        }
    }
    & & \\ \\

    \theutterance \stepcounter{utterance}  
    & & & \multicolumn{2}{p{0.3\linewidth}}{
        \cellcolor[rgb]{0.9,0.9,0.9}{
            \makecell[{{p{\linewidth}}}]{
                \texttt{\tiny{[GM$|$GM]}}
                \texttt{Invalid response format} \\
            }
        }
    }
    & & \\ \\

    \theutterance \stepcounter{utterance}  
    & & & \multicolumn{2}{p{0.3\linewidth}}{
        \cellcolor[rgb]{0.9,0.9,0.9}{
            \makecell[{{p{\linewidth}}}]{
                \texttt{\tiny{[GM$|$GM]}}
                \texttt{Maximum number of penalties exceeded} \\
            }
        }
    }
    & & \\ \\

    \theutterance \stepcounter{utterance}  
    & & & \multicolumn{2}{p{0.3\linewidth}}{
        \cellcolor[rgb]{0.9,0.9,0.9}{
            \makecell[{{p{\linewidth}}}]{
                \texttt{\tiny{[GM$|$GM]}}
                \texttt{Initial states:} \\
\texttt{```} \\
\texttt{Player 1:} \\
\texttt{ 1234567} \\
\texttt{╔═╤══╤══╗} \\
\texttt{║◌│◌◌│◌◌║ 1} \\
\texttt{╟─┤◌◌│◌S║ 2} \\
\texttt{║◌│TD├──╢ 3} \\
\texttt{╟─┘◌◌│O◌║ 4} \\
\texttt{║◌◌PM│◌A║ 5} \\
\texttt{║◌◌┌─┼──╢ 6} \\
\texttt{║◌◌│◌│◌◌║ 7} \\
\texttt{╚══╧═╧══╝} \\
\\ 
\\ 
\texttt{Player 2:} \\
\texttt{ 1234567} \\
\texttt{╔═╤══╤══╗} \\
\texttt{║◌│◌O│T◌║ 1} \\
\texttt{╟─┤◌◌│◌◌║ 2} \\
\texttt{║◌│◌◌├──╢ 3} \\
\texttt{╟─┘◌D│◌◌║ 4} \\
\texttt{║◌SA◌│◌◌║ 5} \\
\texttt{║PM┌─┼──╢ 6} \\
\texttt{║◌◌│◌│◌◌║ 7} \\
\texttt{╚══╧═╧══╝} \\
\\ 
\texttt{```} \\
            }
        }
    }
    & & \\ \\

    \theutterance \stepcounter{utterance}  
    & & & \multicolumn{2}{p{0.3\linewidth}}{
        \cellcolor[rgb]{0.9,0.9,0.9}{
            \makecell[{{p{\linewidth}}}]{
                \texttt{\tiny{[GM$|$GM]}}
                \texttt{End states:} \\
\texttt{```} \\
\texttt{Player 1:} \\
\texttt{ 1234567} \\
\texttt{╔═╤══╤══╗} \\
\texttt{║T│◌◌│◌◌║ 1} \\
\texttt{╟─┤◌◌│◌◌║ 2} \\
\texttt{║◌│PM├──╢ 3} \\
\texttt{╟─┘◌S│◌◌║ 4} \\
\texttt{║◌◌DO│◌A║ 5} \\
\texttt{║◌◌┌─┼──╢ 6} \\
\texttt{║◌◌│◌│◌◌║ 7} \\
\texttt{╚══╧═╧══╝} \\
\\ 
\\ 
\texttt{Player 2:} \\
\texttt{ 1234567} \\
\texttt{╔═╤══╤══╗} \\
\texttt{║◌│◌◌│◌◌║ 1} \\
\texttt{╟─┤◌◌│◌◌║ 2} \\
\texttt{║◌│PM├──╢ 3} \\
\texttt{╟─┘SA│◌◌║ 4} \\
\texttt{║◌◌DO│T◌║ 5} \\
\texttt{║◌◌┌─┼──╢ 6} \\
\texttt{║◌◌│◌│◌◌║ 7} \\
\texttt{╚══╧═╧══╝} \\
\\ 
\texttt{```} \\
            }
        }
    }
    & & \\ \\

    \theutterance \stepcounter{utterance}  
    & & & \multicolumn{2}{p{0.3\linewidth}}{
        \cellcolor[rgb]{0.9,0.9,0.9}{
            \makecell[{{p{\linewidth}}}]{
                \texttt{\tiny{[GM$|$GM]}}
                \texttt{* success: False} \\
\texttt{* lose: False} \\
\texttt{* aborted: True} \\
\texttt{{-}{-}{-}{-}{-}{-}{-}} \\
\texttt{* Shifts: 14.00} \\
\texttt{* Max Shifts: 12.00} \\
\texttt{* Min Shifts: 6.00} \\
\texttt{* End Distance Sum: 10.57} \\
\texttt{* Init Distance Sum: 22.51} \\
\texttt{* Expected Distance Sum: 29.33} \\
\texttt{* Penalties: 18.00} \\
\texttt{* Max Penalties: 16.00} \\
\texttt{* Rounds: 18.00} \\
\texttt{* Max Rounds: 28.00} \\
\texttt{* Object Count: 7.00} \\
            }
        }
    }
    & & \\ \\

\end{supertabular}
}

\end{document}
